% Natural selection & evolution — Biology Upper Sixth — Cours signé OnBuch+
% Official MINESEC syllabus. Compiler avec : tectonic BIO09-natural-selection-evolution.tex
\newcommand{\DOCMATIERE}{Biology}
\newcommand{\DOCNIVEAU}{Upper Sixth}
\newcommand{\DOCMODULE}{Upper Sixth — Biology}
\newcommand{\DOCLECON}{9}
\newcommand{\DOCTITRE}{Natural selection \& evolution}
% OnBuch+ — préambule « cours signés » ENRICHI (compilé avec tectonic / XeLaTeX)
% Système visuel « Pop » d'OnBuch+ : fond crème (jamais blanc), bordures encre
% épaisses, ombres « dures » décalées, titres Archivo Black en capitales.
% Version MATHÉMATIQUES : graphes (pgfplots/tikz), boîtes pédagogiques
% (définition, propriété, méthode, attention, le savais-tu, exercice résolu,
% exercice, TP, bilan), page de garde et signature OnBuch+ ancrée en bas.
%
% Macros attendues dans main.tex AVANT \input{preamble} :
%   \DOCMATIERE  \DOCNIVEAU  \DOCMODULE  \DOCLECON (numéro)  \DOCTITRE
\documentclass[11pt]{article}

\usepackage{fontspec}
\usepackage[english]{babel}
\usepackage[a4paper,margin=2cm,top=3.4cm,bottom=2.8cm,headheight=1.4cm,footskip=1.3cm]{geometry}
\usepackage{amsmath,amssymb,amsfonts}
\usepackage{mathtools} % \xmapsto, \coloneqq... souvent utilisés par les modèles
\usepackage{cancel} % simplifications barrées (bilans de forces)
% \abs{z} (module, valeur absolue) et \norm{u} (norme), utilisés par les modèles
\providecommand{\abs}{}\let\abs\relax\DeclarePairedDelimiter{\abs}{\lvert}{\rvert}
\providecommand{\norm}{}\let\norm\relax\DeclarePairedDelimiter{\norm}{\lVert}{\rVert}
\usepackage[table]{xcolor} % [table] : fournit aussi \rowcolors (alternance de lignes)
\usepackage{pagecolor}
\usepackage{tikz}
\usetikzlibrary{arrows.meta,positioning,calc,shapes.geometric,shapes.misc,shapes.arrows,decorations.pathmorphing,decorations.pathreplacing,decorations.markings,patterns,shadows,mindmap,trees,fit,backgrounds,matrix,intersections,angles,quotes}
\usepackage{pgfplots}
\usepackage[version=4]{mhchem} % équations nucléaires \ce{^{235}_{92}U}
\usepackage[european,siunitx]{circuitikz} % schémas électriques
\pgfplotsset{compat=1.18}
\usepgfplotslibrary{fillbetween,groupplots} % aires entre courbes (calcul intégral)
\usepackage{enumitem}
\usepackage{fancyhdr}
% [most] charge toutes les bibliothèques tcolorbox (listings, minted, raster,
% documentation...) et gonfle inutilement la mémoire TeX sur un document long
% (~300 boîtes) : on ne charge que ce qui est réellement utilisé.
\usepackage[skins,breakable]{tcolorbox}
\usepackage{graphicx}
\usepackage{colortbl}
\usepackage{booktabs}
\usepackage{tabularx}
\usepackage{diagbox} % case d'en-tête diagonale des tableaux à double entrée
% Colonnes extensibles centrée / gauche / droite (C, L, R), souvent utilisées par les modèles
\newcolumntype{C}{>{\centering\arraybackslash}X}
\newcolumntype{L}{>{\raggedright\arraybackslash}X}
\newcolumntype{R}{>{\raggedleft\arraybackslash}X}
\usepackage{array}
\usepackage{multicol}
\usepackage{siunitx}
% parse-numbers=false : accepte aussi \SI{2,5 \times 10^{-3}}{...}
\sisetup{locale=UK,output-decimal-marker={.},group-minimum-digits=4,parse-numbers=false}
\DeclareSIUnit\ohm{\ensuremath{\symup{\Omega}}} % U+2126 absent de la police
\DeclareSIUnit\division{div} % réglages d'oscilloscope (ms/div, V/div)
\DeclareSIUnit\tr{tr} % tours (vitesse de rotation, ex. \SI{1500}{\tr\per\minute})
\usepackage{qrcode}
% \degree : commande fréquemment utilisée par les modèles pour un angle ou une
% température en dehors de \SI{}{} (non fournie par les paquets chargés).
\providecommand{\degree}{^{\circ}}
% \qed (fin de démonstration), utilisé par les modèles sans amsthm
\providecommand{\qed}{\hfill$\square$}
% \barre : double barre des valeurs interdites dans un tableau de variations (tkz-tab)
\providecommand{\barre}{\ensuremath{\|}}
% environnement proof (amsthm non chargé) utilisé par les modèles
\ifdefined\proof\else\newenvironment{proof}[1][Proof]{\par\noindent\textit{#1.}\ }{\qed\par}\fi
\usepackage{lastpage}
\usepackage{hyperref}
\hypersetup{hidelinks}

% ============================================================
% Palette « Pop » OnBuch+ — mêmes hex que la classe P (plus_ui.dart)
% ============================================================
\definecolor{poporange}{HTML}{F59321}
\definecolor{popdark}{HTML}{E07A0C}
\definecolor{popcream}{HTML}{FFF6E8}
\definecolor{popink}{HTML}{1C1714}
\definecolor{poppurple}{HTML}{7A5AE0}
\definecolor{popgreen}{HTML}{2E9E6A}
\definecolor{poppink}{HTML}{E0457B}
\definecolor{popblue}{HTML}{2F7DE1}
\definecolor{popgold}{HTML}{C98A12}
\definecolor{popmuted}{HTML}{8A7F76}
\definecolor{popline}{HTML}{EADFD2}
\definecolor{popgray}{HTML}{8A7F76}
\colorlet{popgrey}{popgray}
% Alias tolérants (le modèle nomme parfois les couleurs différemment)
\colorlet{popwhite}{white}
\colorlet{popblack}{popink}
\colorlet{popred}{poppink}
\colorlet{popyellow}{popgold}
% Teintes claires pour fonds de tableaux / zones
\colorlet{popblueL}{popblue!10!white}
\colorlet{poporangeL}{poporange!14!white}
\colorlet{popgreenL}{popgreen!12!white}
\colorlet{poppurpleL}{poppurple!12!white}
\colorlet{poppinkL}{poppink!12!white}
\colorlet{popgoldL}{popgold!14!white}

\pagecolor{popcream}

% ============================================================
% Typographie
% ============================================================
\defaultfontfeatures{Ligatures=TeX}
\setmainfont{PlusJakartaSans-Medium.ttf}[Path=../../../fonts/,BoldFont=PlusJakartaSans-Bold.ttf]
% Maths en sans-serif (Fira Math) avec chiffres et lettres droites de Plus
% Jakarta, pour que les formules se fondent dans le texte.
\usepackage[math-style=ISO,bold-style=ISO]{unicode-math}
\setmathfont{FiraMath-Regular.otf}[Path=../../../fonts/]
\setmathfont{PlusJakartaSans-Medium.ttf}[Path=../../../fonts/,range={up/num,up/{Latin,latin},\mathup}]
% (icomma retiré : décimales avec point en anglais)
% Caractères mathématiques Unicode absents des polices de texte (Plus Jakarta,
% Archivo Black) : rendus par leur équivalent mathématique, en texte comme en
% formule — sinon ils s'affichent en carré vide (ex. « Arithmétique dans ℤ »).
\usepackage{newunicodechar}
\newunicodechar{ℝ}{\ensuremath{\mathbb{R}}}
\newunicodechar{ℤ}{\ensuremath{\mathbb{Z}}}
\newunicodechar{ℂ}{\ensuremath{\mathbb{C}}}
\newunicodechar{ℕ}{\ensuremath{\mathbb{N}}}
\newunicodechar{ℚ}{\ensuremath{\mathbb{Q}}}
\newunicodechar{𝔻}{\ensuremath{\mathbb{D}}}
\newunicodechar{α}{\ensuremath{\alpha}}
\newunicodechar{β}{\ensuremath{\beta}}
\newunicodechar{γ}{\ensuremath{\gamma}}
\newunicodechar{δ}{\ensuremath{\delta}}
\newunicodechar{ε}{\ensuremath{\varepsilon}}
\newunicodechar{θ}{\ensuremath{\theta}}
\newunicodechar{λ}{\ensuremath{\lambda}}
\newunicodechar{μ}{\ensuremath{\mu}}
\newunicodechar{σ}{\ensuremath{\sigma}}
\newunicodechar{τ}{\ensuremath{\tau}}
\newunicodechar{φ}{\ensuremath{\varphi}}
\newunicodechar{ψ}{\ensuremath{\psi}}
\newunicodechar{ω}{\ensuremath{\omega}}
\newunicodechar{Σ}{\ensuremath{\Sigma}}
% majuscules produites par \MakeUppercase dans les titres (α-aminés -> Α)
\newunicodechar{Α}{\ensuremath{\alpha}}
\newunicodechar{Β}{\ensuremath{\beta}}
\newunicodechar{Γ}{\ensuremath{\Gamma}}
\newunicodechar{Δ}{\ensuremath{\Delta}}
\newunicodechar{Ε}{\ensuremath{\varepsilon}}
\newunicodechar{Θ}{\ensuremath{\Theta}}
\newunicodechar{Λ}{\ensuremath{\Lambda}}
\newunicodechar{Μ}{\ensuremath{\mu}}
\newunicodechar{Τ}{\ensuremath{\tau}}
\newunicodechar{Φ}{\ensuremath{\Phi}}
\newunicodechar{Ψ}{\ensuremath{\Psi}}
\newunicodechar{Ω}{\ensuremath{\Omega}}
\newunicodechar{↦}{\ensuremath{\mapsto}}
\newunicodechar{∈}{\ensuremath{\in}}
\newunicodechar{∉}{\ensuremath{\notin}}
\newunicodechar{∧}{\ensuremath{\wedge}}
\newunicodechar{⇔}{\ensuremath{\Leftrightarrow}}
\newunicodechar{⟺}{\ensuremath{\Longleftrightarrow}}
\newunicodechar{⇒}{\ensuremath{\Rightarrow}}
\newunicodechar{⟹}{\ensuremath{\Longrightarrow}}
\newunicodechar{∘}{\ensuremath{\circ}}
\newunicodechar{∪}{\ensuremath{\cup}}
\newunicodechar{∩}{\ensuremath{\cap}}
\newunicodechar{≡}{\ensuremath{\equiv}}
\newunicodechar{⊂}{\ensuremath{\subset}}
\newunicodechar{⊥}{\ensuremath{\perp}}
\newunicodechar{∃}{\ensuremath{\exists}}
\newunicodechar{∀}{\ensuremath{\forall}}
\newunicodechar{∛}{\ensuremath{\sqrt[3]{\,}}}
\newunicodechar{∜}{\ensuremath{\sqrt[4]{\,}}}
\newunicodechar{⃗}{\ensuremath{{}^{\to}}}
\newunicodechar{ⁿ}{\ifmmode^{n}\else\textsuperscript{\ensuremath{n}}\fi}
\newunicodechar{ˣ}{\ifmmode^{x}\else\textsuperscript{\ensuremath{x}}\fi}
\newunicodechar{ᵇ}{\ifmmode^{b}\else\textsuperscript{\ensuremath{b}}\fi}
\newunicodechar{ʳ}{\ifmmode^{r}\else\textsuperscript{\ensuremath{r}}\fi}
\newunicodechar{ᵘ}{\ifmmode^{u}\else\textsuperscript{\ensuremath{u}}\fi}
\newunicodechar{ʸ}{\ifmmode^{y}\else\textsuperscript{\ensuremath{y}}\fi}
\newunicodechar{ᵃ}{\ifmmode^{a}\else\textsuperscript{\ensuremath{a}}\fi}
\newunicodechar{ᵉ}{\ifmmode^{e}\else\textsuperscript{\ensuremath{e}}\fi}
\newunicodechar{ᵏ}{\ifmmode^{k}\else\textsuperscript{\ensuremath{k}}\fi}
\newunicodechar{ᵖ}{\ifmmode^{p}\else\textsuperscript{\ensuremath{p}}\fi}
\newunicodechar{ᴺ}{\ifmmode^{N}\else\textsuperscript{\ensuremath{N}}\fi}
\newunicodechar{ᶜ}{\ifmmode^{c}\else\textsuperscript{\ensuremath{c}}\fi}
\newunicodechar{ʰ}{\ifmmode^{h}\else\textsuperscript{\ensuremath{h}}\fi}
\newunicodechar{ᵠ}{\ifmmode^{q}\else\textsuperscript{\ensuremath{q}}\fi}
\newunicodechar{ₙ}{\ifmmode_{n}\else\textsubscript{\ensuremath{n}}\fi}
\newunicodechar{ₐ}{\ifmmode_{a}\else\textsubscript{\ensuremath{a}}\fi}
\newunicodechar{ᵢ}{\ifmmode_{i}\else\textsubscript{\ensuremath{i}}\fi}
\newunicodechar{ᵣ}{\ifmmode_{r}\else\textsubscript{\ensuremath{r}}\fi}
\newunicodechar{ₖ}{\ifmmode_{k}\else\textsubscript{\ensuremath{k}}\fi}
\newunicodechar{ᵦ}{\ifmmode_{\beta}\else\textsubscript{\ensuremath{\beta}}\fi}
\newunicodechar{ₓ}{\ifmmode_{x}\else\textsubscript{\ensuremath{x}}\fi}
\newfontfamily\popdisplay{ArchivoBlack-Regular.ttf}[Path=../../../fonts/]
\newfontfamily\poplogo{SpaceGrotesk-Bold.ttf}[Path=../../../fonts/]
\newfontfamily\popsemi{PlusJakartaSans-SemiBold.ttf}[Path=../../../fonts/]
\newfontfamily\popxbold{PlusJakartaSans-ExtraBold.ttf}[Path=../../../fonts/]
\newfontfamily\popxboldls{PlusJakartaSans-ExtraBold.ttf}[Path=../../../fonts/,LetterSpace=3.0]

\setlist[enumerate]{leftmargin=1.7em,itemsep=5pt,topsep=4pt}
\setlist[itemize]{leftmargin=1.7em,itemsep=4pt,topsep=4pt,label={\color{poporange}\textbullet}}
\setlength{\parindent}{0pt}
\setlength{\parskip}{5pt}
\renewcommand{\arraystretch}{1.25}

% pgfplots : style Pop par défaut
\pgfplotsset{
  every axis/.append style={axis line style={popink,line width=1pt},tick style={popink},
    grid style={popline,line width=0.5pt},label style={font=\small},tick label style={font=\footnotesize}},
  popcurve/.style={poporange,line width=1.8pt,no markers,smooth},
}
% Même style disponible dans un \draw TikZ simple (hors pgfplots)
\tikzset{popcurve/.style={poporange,line width=1.8pt}}
\tikzset{popgrid/.style={popline,very thin}}
\colorlet{popcurve}{poporange} % le nom du style est parfois utilisé comme couleur

% ============================================================
% Pill / squiggle
% ============================================================
\newcommand{\poppill}[3]{%
  \tikz[baseline=-0.55ex]\node[draw=popink,line width=1.2pt,rounded corners=2.6pt,%
    fill=#2,inner xsep=6.5pt,inner ysep=3pt]{%
    \popxboldls\fontsize{7}{7}\selectfont\color{#3}\MakeUppercase{#1}};%
}
\newcommand{\popsquiggle}[1]{%
  \tikz[baseline=-0.4ex]{
    \draw[#1,line width=2.6pt,line cap=round]
      plot [smooth] coordinates {(0,0.09) (0.34,0.01) (0.68,0.21) (1.02,0.01) (1.36,0.10)};
  }%
}
% Mot-clé surligné (marqueur orange sous le texte)
\newcommand{\cle}[1]{{\popxbold\color{popdark}#1}}

% ============================================================
% En-tête / pied
% ============================================================
\pagestyle{fancy}
\fancyhf{}
\renewcommand{\headrulewidth}{0pt}
\renewcommand{\footrulewidth}{0pt}
\lhead{\raisebox{-0.15ex}{\poplogo\fontsize{13}{13}\selectfont\color{popink}OnBuch\color{poporange}+}}
\rhead{\poppill{\DOCMATIERE\ \textbullet\ \DOCNIVEAU\ \textbullet\ Chapter \DOCLECON}{white}{popink}}
\lfoot{\popxbold\fontsize{7.6}{7.6}\selectfont\color{popmuted}\MakeUppercase{Signed course OnBuch+}\ \textbullet\ {\popsemi\DOCTITRE}}
\rfoot{\tikz[baseline=-0.6ex]\node[rounded corners=8.5pt,fill=popink,minimum height=17pt,minimum width=17pt,inner xsep=5pt,inner sep=0]{\popxbold\fontsize{8}{8}\selectfont\color{popcream}\thepage\,/\,\pageref*{LastPage}};}

% ============================================================
% Titres
% ============================================================
\newcommand{\chaptitre}[2]{%
  \begin{tcolorbox}[enhanced,colback=poporange,colframe=popink,boxrule=2.6pt,arc=11pt,
      drop shadow=popink,left=15pt,right=15pt,top=12pt,bottom=11pt,before skip=3pt,after skip=18pt]
    \poppill{#2}{popink}{popcream}\par\vspace{9pt}
    {\raggedright\hyphenpenalty=10000\exhyphenpenalty=10000\popdisplay\fontsize{22}{24}\selectfont\color{popcream}\MakeUppercase{#1}\par}
  \end{tcolorbox}
}
% Section numérotée : pastille orange avec le numéro + titre Archivo
\newcounter{popsec}
\newcommand{\coursec}[1]{%
  \stepcounter{popsec}\setcounter{popsub}{0}%
  \par\vspace{16pt}\noindent
  \begin{minipage}{\linewidth}
  \tikz[baseline=-0.7ex]\node[draw=popink,line width=1.6pt,fill=poporange,rounded corners=4pt,minimum size=22pt,inner sep=0]{\popdisplay\fontsize{12}{12}\selectfont\color{popcream}\thepopsec};%
  \hspace{8pt}{\popdisplay\fontsize{15}{17}\selectfont\color{popink}\MakeUppercase{#1}}\par
  \vspace{4pt}\noindent{\color{poporange}\rule{36pt}{3.6pt}}
  \end{minipage}\par\nopagebreak\vspace{8pt}
}
\newcounter{popsub}
\newcommand{\courssub}[1]{%
  \stepcounter{popsub}%
  \par\vspace{9pt}\noindent{\popxbold\fontsize{11.5}{13}\selectfont\color{popink}{\color{poporange}\thepopsec.\thepopsub}\hspace{6pt}#1}\par\nopagebreak\vspace{4pt}
}

% ============================================================
% Boîtes pédagogiques — même recette : fond blanc, bordure épaisse colorée,
% ombre dure encre, badge plein. Titre optionnel [#1].
% ============================================================
\tcbset{popbase/.style={breakable,enhanced,colback=white,boxrule=2.3pt,arc=9pt,
  shadow={3pt}{-3pt}{0pt}{popink},left=14pt,right=14pt,top=11pt,bottom=11pt,before skip=11pt,after skip=15pt,
  fontupper=\popsemi\fontsize{10.6}{14.5}\selectfont\color{popink},
  fontlower=\popsemi\fontsize{10.6}{14.5}\selectfont\color{popink}}}
% Coche dessinée (le glyphe ✓ n'existe pas dans les polices du document)
\newcommand{\popcheck}[1]{\tikz[baseline=-0.5ex]\draw[#1,line width=1.6pt,line cap=round,line join=round](-0.13,0.0)--(-0.04,-0.09)--(0.13,0.1);}
\providecommand{\coche}{\popcheck{popgreen}}
\providecommand{\resultat}[1]{\par\smallskip\noindent\fcolorbox{popgreen}{popgreenL}{\parbox{\dimexpr\linewidth-2\fboxsep-2\fboxrule\relax}{\textbf{Result.} #1}}\par\smallskip}
\newcommand{\popboxhead}[3]{\poppill{#1}{#2}{white}\ifx\relax#3\relax\else\hspace{6pt}{\popxbold\fontsize{10.6}{12}\selectfont\color{popink}#3}\fi\par\vspace{7pt}}

\newtcolorbox{aretenir}[1][]{popbase,colframe=popgreen,before upper={\popboxhead{Key points}{popgreen}{#1}}}
\newtcolorbox{exemplebox}[1][]{popbase,colframe=poporange,before upper={\popboxhead{Exemple}{poporange}{#1}}}
\newtcolorbox{definition}[1][]{popbase,colframe=popblue,before upper={\popboxhead{Definition}{popblue}{#1}}}
\newtcolorbox{propriete}[1][]{popbase,colframe=poppurple,before upper={\popboxhead{Property}{poppurple}{#1}}}
\newtcolorbox{methode}[1][]{popbase,colframe=popgold,before upper={\popboxhead{Method}{popgold}{#1}}}
\newtcolorbox{attention}[1][]{popbase,colframe=poppink,before upper={\popboxhead{Warning}{poppink}{#1}}}
\newtcolorbox{experience}[1][]{popbase,colframe=popblue,colback=popblueL,before upper={\popboxhead{Experiment}{popblue}{#1}}}
% « Le savais-tu ? » : carte violette pleine (contexte, culture, Cameroun)
\newtcolorbox{savaistu}[1][]{popbase,colback=poppurple,colframe=popink,
  fontupper=\popsemi\fontsize{10.4}{14.2}\selectfont\color{white},
  before upper={\poppill{Did you know?}{white}{poppurple}\ifx\relax#1\relax\else\hspace{6pt}{\popxbold\color{white}#1}\fi\par\vspace{7pt}}}
% Exercice résolu : énoncé en haut, \tcblower puis la solution sur fond crème
\newcounter{popexo}\newcounter{popexr}
\newtcolorbox{exoresolu}[1][]{popbase,colframe=poporange,colbacklower=popcream,
  segmentation style={popink,line width=1.2pt,dashed},
  before upper={\stepcounter{popexr}\popboxhead{Worked example \thepopexr}{poporange}{#1}},
  before lower={\poppill{Solution}{popink}{popcream}\par\vspace{6pt}}}
% Exercice (non résolu en place) — la correction est dans la partie Corrigés
\newtcolorbox{exercice}[1][]{popbase,colframe=popink,
  before upper={\stepcounter{popexo}\popboxhead{Exercise \thepopexo}{popink}{#1}}}
% Corrigé
\newtcolorbox{corrige}[1][]{popbase,colframe=popgreen,colback=popgreenL,
  before upper={\popboxhead{Solution}{popgreen}{#1}}}
% Objectifs / prérequis / bilan
\newtcolorbox{objectifs}{popbase,colframe=popink,colback=poporangeL,
  before upper={\popboxhead{Chapter objectives}{popink}{}}}
\newtcolorbox{prerequis}{popbase,colframe=popink,colback=popblueL,
  before upper={\popboxhead{Prerequisites}{popblue}{}}}
\newtcolorbox{bilan}{popbase,colframe=popink,colback=white,boxrule=2.8pt,
  before upper={\popboxhead{Fiche bilan}{poporange}{L'essentiel à connaître pour l'examen}}}
% Figure encadrée avec légende
\newtcolorbox{popfigure}[1][]{enhanced,breakable=false,colback=white,colframe=popline,boxrule=1.4pt,arc=8pt,
  left=10pt,right=10pt,top=10pt,bottom=8pt,before skip=10pt,after skip=12pt,halign=center,
  lower separated=false,halign lower=center,
  fontlower=\popsemi\fontsize{9}{11.5}\selectfont\color{popmuted},
  #1}
\newcommand{\legende}[1]{\par\vspace{6pt}{\popsemi\fontsize{9}{11.5}\selectfont\color{popmuted}#1}}

% ============================================================
% Signature OnBuch+
% ============================================================
\newcommand{\poplogoinline}[1]{{\poplogo\fontsize{#1}{#1}\selectfont\color{popink}OnBuch\color{poporange}+}}
\newcommand{\signatureonbuch}{%
  \begin{tcolorbox}[enhanced,colback=popink,colframe=popink,boxrule=2.2pt,arc=10pt,
      drop shadow=poporange,left=14pt,right=14pt,top=10pt,bottom=10pt,before skip=0pt,after skip=0pt]
    \tikz[baseline=-0.7ex]\node[circle,fill=popgreen,draw=popcream,line width=1.3pt,minimum size=22pt,inner sep=0]{\popcheck{white}};%
    \hspace{9pt}\begin{minipage}[c]{0.62\linewidth}
      {\popxbold\fontsize{12}{13}\selectfont\color{popcream}Signed\ \textbullet\ {\poplogo OnBuch\color{poporange}+}}\par\vspace{2pt}
      {\raggedright\popsemi\fontsize{8}{10.5}\selectfont\color{popline}\MakeUppercase{OnBuch+ teaching team}\\\MakeUppercase{Follows the official MINESEC syllabus}\par}
    \end{minipage}\hfill
    {\poplogo\fontsize{22}{22}\selectfont\color{popcream}OnBuch\color{poporange}+}
  \end{tcolorbox}%
}

% ============================================================
% Page de garde
% #1 titre · #2 matière · #3 niveau · #4 module · #5 n° leçon · #6 durée ·
% #7 description / objectifs (texte ou itemize)
% ============================================================
\newcommand{\pagegarde}[7]{%
  \thispagestyle{empty}
  \newgeometry{a4paper,margin=1.7cm,top=1.6cm,bottom=1.4cm}%
  \begin{tikzpicture}[remember picture,overlay]
    \fill[poporange!18] ([xshift=-3.2cm,yshift=-3.6cm]current page.north east) circle (4.6cm);
    \fill[poppurple!14] ([xshift=2.4cm,yshift=7.6cm]current page.south west) circle (3.8cm);
  \end{tikzpicture}%
  {\poplogo\fontsize{30}{30}\selectfont\color{popink}OnBuch\color{poporange}+}\hfill
  \raisebox{0.6ex}{\poppill{Signed course \textbullet\ Premium}{popink}{popcream}}\par
  \vspace{1pt}\popsquiggle{poppurple}\par
  \vspace{8pt}
  \begin{tcolorbox}[enhanced,colback=poporange,colframe=popink,boxrule=2.8pt,arc=13pt,
      drop shadow=popink,left=18pt,right=18pt,top=12pt,bottom=13pt,after skip=10pt]
    \poppill{#2 \textbullet\ #3}{popink}{popcream}\hspace{5pt}\poppill{#4}{white}{popink}\par\vspace{8pt}
    {\popdisplay\fontsize{14}{14}\selectfont\color{popink}CHAPTER #5}\par\vspace{4pt}
    {\raggedright\hyphenpenalty=10000\exhyphenpenalty=10000\popdisplay\fontsize{25}{27}\selectfont\color{popcream}\MakeUppercase{#1}\par}
    \vspace{8pt}
    {\popxbold\fontsize{9.5}{9.5}\selectfont\color{popink}\MakeUppercase{Suggested time: #6}}
  \end{tcolorbox}
  \begin{tcolorbox}[enhanced,colback=white,colframe=popink,boxrule=2.2pt,arc=10pt,
      drop shadow=popink,left=16pt,right=16pt,top=10pt,bottom=10pt,after skip=8pt]
    \poppill{In this chapter}{poppurple}{white}\par\vspace{5pt}
    \popsemi\fontsize{9.3}{12.6}\selectfont\color{popink}#7
  \end{tcolorbox}
  \noindent\begin{minipage}[t]{0.487\linewidth}
    \begin{tcolorbox}[enhanced,colback=popgreen,colframe=popink,boxrule=2.2pt,arc=10pt,
        drop shadow=popink,left=11pt,right=11pt,top=10pt,bottom=10pt,height=3.1cm,valign=center]
      \tcbox[colback=white,colframe=popink,boxrule=1.3pt,arc=5pt,boxsep=0pt,nobeforeafter,
        left=3pt,right=3pt,top=3pt,bottom=3pt,valign=center]{\qrcode[height=1.9cm]{https://play.google.com/store/apps/details?id=com.luvvix.onbuch&hl=en}}%
      \hspace{8pt}\begin{minipage}[c]{0.5\linewidth}
        {\popdisplay\fontsize{11}{12.5}\selectfont\color{white}\MakeUppercase{Download OnBuch}}\par\vspace{4pt}
        {\raggedright\popsemi\fontsize{8.4}{11}\selectfont\color{white}Free \textbullet\ Google Play\\AI tutor, past papers, results\par}
      \end{minipage}
    \end{tcolorbox}
  \end{minipage}\hfill
  \begin{minipage}[t]{0.487\linewidth}
    \begin{tcolorbox}[enhanced,colback=poppurple,colframe=popink,boxrule=2.2pt,arc=10pt,
        drop shadow=popink,left=11pt,right=11pt,top=10pt,bottom=10pt,height=3.1cm,valign=center]
      \tikz[baseline=-0.7ex]\node[circle,fill=white,draw=popink,line width=1.3pt,minimum size=1.6cm,inner sep=0]{\popdisplay\fontsize{24}{24}\selectfont\color{poppurple}+};%
      \hspace{8pt}\begin{minipage}[c]{0.6\linewidth}
        {\popdisplay\fontsize{11}{12.5}\selectfont\color{white}\MakeUppercase{Go OnBuch+}}\par\vspace{4pt}
        {\raggedright\popsemi\fontsize{8.4}{11}\selectfont\color{white}All signed courses, unlimited AI tutor, PDF sheets — OnBuch+ tab in the app\par}
      \end{minipage}
    \end{tcolorbox}
  \end{minipage}
  \vspace{8pt}
  \popcontenu
  \vfill
  \signatureonbuch
  \restoregeometry
  \newpage
}

% Rangée « Ce cours contient » (page de garde)
\newcommand{\poptile}[3]{% #1 couleur #2 gros texte #3 légende
  \begin{tcolorbox}[enhanced,colback=#1,colframe=popink,boxrule=2pt,arc=9pt,shadow={2.5pt}{-2.5pt}{0pt}{popink},
      left=6pt,right=6pt,top=9pt,bottom=9pt,halign=center,valign=center,height=1.75cm,width=\linewidth,nobeforeafter]
    {\popdisplay\fontsize{12.5}{13}\selectfont\color{white}\MakeUppercase{#2}}\par\vspace{3pt}
    {\centering\popsemi\fontsize{7.8}{9.5}\selectfont\color{white}#3\par}
  \end{tcolorbox}}
\newcommand{\popcontenu}{%
  \par\noindent{\popxboldls\fontsize{7.5}{7.5}\selectfont\color{popmuted}THIS SIGNED COURSE CONTAINS}\par\vspace{7pt}
  \noindent\begin{minipage}[t]{0.235\linewidth}\poptile{popblue}{Course}{complete, illustrated, step by step}\end{minipage}\hfill
  \begin{minipage}[t]{0.235\linewidth}\poptile{poporange}{Methods}{and worked examples}\end{minipage}\hfill
  \begin{minipage}[t]{0.235\linewidth}\poptile{poppink}{Exercises}{exam style + solutions}\end{minipage}\hfill
  \begin{minipage}[t]{0.235\linewidth}\poptile{popgreen}{Summary sheet}{the essentials to remember}\end{minipage}\par}

% Sommaire
\newtcolorbox{sommaire}{enhanced,colback=white,colframe=popink,boxrule=2.2pt,arc=10pt,
  drop shadow=popink,left=18pt,right=18pt,top=13pt,bottom=13pt,before skip=6pt,after skip=20pt,
  before upper={\poppill{Contents}{poppurple}{white}\par\vspace{9pt}\popsemi\fontsize{10.6}{14.5}\selectfont\color{popink}}}

% Signature de fin de cours — poussée tout en bas de la dernière page
\newcommand{\findecours}{%
  \par\vfill\nopagebreak
  \begin{center}\popsquiggle{poporange}\end{center}\vspace{4pt}
  \signatureonbuch
}
\AtBeginDocument{\def\llbracket{\mathopen{[\mkern-3.5mu[}}\def\rrbracket{\mathclose{]\mkern-3.5mu]}}} % intervalles d'entiers (glyphe ⟦ absent de la police)
% Glyphes absents de Fira Math : redéfinis à partir de glyphes disponibles
\AtBeginDocument{%
  \def\setminus{\mathbin{\backslash}}\let\smallsetminus\setminus
  \def\vdots{\mathord{\vbox{\baselineskip3.2pt\lineskiplimit0pt\kern4pt\hbox{.}\hbox{.}\hbox{.}}}}%
  \def\ddots{\mathinner{\mkern1mu\raise6.5pt\hbox{.}\mkern2mu\raise3.6pt\hbox{.}\mkern2mu\raise0.7pt\hbox{.}\mkern1mu}}%
  \def\triangle{\mathord{\scalebox{0.9}{$\symup{\Delta}$}}}%
  \def\nabla{\mathord{\raisebox{\depth}{\scalebox{1}[-1]{$\symup{\Delta}$}}}}%
  \def\checkmark{\popcheck{popdark}}%
  \def\star{\mathbin{\ast}}\def\bigstar{\mathord{\ast}}%
  \def\diamond{\mathbin{\scalebox{0.6}{$\Diamond$}}}%
  \def\bigcup{\mathop{\mathchoice{\raisebox{-0.3em}{\scalebox{1.7}{$\cup$}}}{\scalebox{1.25}{$\cup$}}{\cup}{\cup}}\displaylimits}%
  \def\bigcap{\mathop{\mathchoice{\raisebox{-0.3em}{\scalebox{1.7}{$\cap$}}}{\scalebox{1.25}{$\cap$}}{\cap}{\cap}}\displaylimits}%
  \def\lesseqqgtr{\mathrel{\lessgtr}}\def\bigoplus{\mathop{\oplus}\displaylimits}%
}
\providecommand{\ssi}{if and only if} % abréviation fréquente
\AtBeginDocument{\providecommand{\wideparen}{\overparen}} % arc de cercle

%% FIGFIX-BEGIN v4 — correctifs globaux des figures (docs/onbuch-generation/tools/figfix)
\usepackage{adjustbox}\usepackage{etoolbox}
% toute figure TikZ/pgfplots plus large que la ligne est réduite à la largeur disponible
\BeforeBeginEnvironment{tikzpicture}{\begin{adjustbox}{max width=\linewidth}}
\AfterEndEnvironment{tikzpicture}{\end{adjustbox}}
% légendes pgfplots lisibles par-dessus les courbes
\pgfplotsset{every axis/.append style={legend style={fill=white,fill opacity=0.92,text opacity=1,draw=black!15,inner sep=3pt}}}
% étiquettes d'arêtes (syntaxe edge["..."]) avec fond blanc
\tikzset{every edge quotes/.append style={fill=white,inner sep=1.2pt}}
%% FIGFIX-END
\begin{document}

\pagegarde{Natural selection \& evolution}{Biology}{Upper Sixth}{Upper Sixth — Biology}{9}{6 periods}{This chapter explores how natural selection shapes populations and leads to the formation of new species, the evidence for organic evolution, and the major theories proposed to explain it. It prepares Upper Sixth candidates for the GCE Advanced Level theory and practical examinations in Biology.
\vspace{4pt}
\begin{multicols}{2}\small
\begin{itemize}
\item By the end of this chapter I can define evolution, natural selection, adaptation, fitness and speciation using precise biological vocabulary.
\item By the end of this chapter I can explain the mechanism of natural selection using variation, overproduction, competition and differential survival.
\item By the end of this chapter I can apply the Hardy-Weinberg principle to calculate allele and genotype frequencies in a population.
\item By the end of this chapter I can distinguish directional, stabilising and disruptive selection with real examples.
\item By the end of this chapter I can explain how isolating mechanisms lead to speciation (allopatric and sympatric).
\item By the end of this chapter I can describe and evaluate the main lines of evidence for evolution: fossils, comparative anatomy, embryology, molecular biology and biogeography.
\end{itemize}\end{multicols}}

\begin{sommaire}
\begin{multicols}{2}
\begin{enumerate}
\item Section 1: Variation, Populations and the Genetic Basis of Evolution
\item Section 2: Natural Selection — Mechanism and Types
\item Section 3: Speciation and Isolating Mechanisms
\item Section 4: Evidence of Organic Evolution
\item Section 5: Theories of Evolution and GCE Revision
\item Integration activity
\item Methods and worked examples
\item Exercises
\item Solutions to the exercises
\item Summary sheet
\end{enumerate}
\end{multicols}
\end{sommaire}

\begin{objectifs}
\begin{itemize}
\item By the end of this chapter I can define evolution, natural selection, adaptation, fitness and speciation using precise biological vocabulary.
\item By the end of this chapter I can explain the mechanism of natural selection using variation, overproduction, competition and differential survival.
\item By the end of this chapter I can apply the Hardy-Weinberg principle to calculate allele and genotype frequencies in a population.
\item By the end of this chapter I can distinguish directional, stabilising and disruptive selection with real examples.
\item By the end of this chapter I can explain how isolating mechanisms lead to speciation (allopatric and sympatric).
\item By the end of this chapter I can describe and evaluate the main lines of evidence for evolution: fossils, comparative anatomy, embryology, molecular biology and biogeography.
\item By the end of this chapter I can compare the theories of Lamarck, Darwin and the modern synthetic theory (neo-Darwinism).
\item By the end of this chapter I can interpret phylogenetic trees and data on allele frequency changes in exam-style questions.
\item By the end of this chapter I can answer structured and essay GCE questions on evolution within the allotted time.
\end{itemize}
\end{objectifs}

\begin{prerequis}
\begin{itemize}
\item Cell structure, chromosomes, genes and alleles (Lower Sixth: cell biology and genetics)
\item Mendelian inheritance, monohybrid and dihybrid crosses, genotype and phenotype
\item Variation within populations: continuous and discontinuous variation, mutation as a source of variation
\item DNA structure, protein synthesis and gene mutation
\item Classification of living organisms and the concept of a species
\end{itemize}
\end{prerequis}

\newpage

\coursec{Section 1: Variation, Populations and the Genetic Basis of Evolution}

In the markets of Yaoundé and Bamenda, farmers complain that the same pesticides that killed mosquitoes and crop pests ten years ago barely work today. Meanwhile, doctors in Douala warn that some malaria parasites no longer respond to drugs that once cured patients easily. Why do living populations seem to ``fight back'' against our chemicals? The answer lies in one of the most powerful ideas in biology: \cle{natural selection} drives populations to change over generations. In this chapter, you will discover how evolution works, what evidence supports it, and how new species arise --- knowledge that is heavily examined at the GCE Advanced Level (Lessons 93--95).

Before we can understand how populations change, we must first understand what a population is in genetic terms, where its variation comes from, and how we can measure that variation mathematically. That is the goal of this first section.

\courssub{Populations, Gene Pools and Allele Frequencies}

\textbf{From individuals to populations.}\par
Evolution does not happen to individuals: a single mosquito cannot ``evolve'' resistance during its lifetime. What changes is the \emph{population} --- the proportion of resistant individuals rises from one generation to the next. This is why the \cle{population} is the true unit of evolution.

\begin{definition}[Population, gene pool and allele frequency]
A \cle{population} is a group of organisms of the \textbf{same species} living in the \textbf{same area} at the \textbf{same time}, capable of interbreeding and producing fertile offspring.

The \cle{gene pool} of a population is the \textbf{total collection of all the alleles} of all the genes present in that population at a given time.

The \cle{allele frequency} of a given allele is the proportion of that allele among all copies of the gene in the population:
\[
\text{allele frequency}=\frac{\text{number of copies of the allele}}{\text{total number of copies of the gene in the population}}
\]
The \cle{genotype frequency} is the proportion of individuals in the population carrying a particular genotype (e.g.\ $AA$, $Aa$ or $aa$).
\end{definition}

\begin{exemplebox}[Counting alleles in a small population]
Imagine a population of 50 plants in which a flower-colour gene has two alleles, $R$ (red) and $r$ (white). Suppose there are 20 plants of genotype $RR$, 20 of genotype $Rr$ and 10 of genotype $rr$. Each plant carries \textbf{two} copies of the gene, so the gene pool contains $50 \times 2 = 100$ alleles.
\begin{itemize}
\item Copies of $R$: $(20 \times 2) + 20 = 60$, so frequency of $R = \dfrac{60}{100} = 0.6$.
\item Copies of $r$: $(10 \times 2) + 20 = 40$, so frequency of $r = \dfrac{40}{100} = 0.4$.
\end{itemize}
Notice that $0.6 + 0.4 = 1$: the frequencies of all alleles of a gene always add up to 1.
\end{exemplebox}

\begin{popfigure}
\begin{tikzpicture}[scale=0.95]
% Generation 1
\draw[fill=popblueL, draw=popblue, thick] (0,0) ellipse (1.5 and 1.0);
\node[font=\bfseries\small, text=popdark] at (0,1.35) {Generation 1};
\foreach \x/\y in {-0.7/0.3, -0.2/0.5, 0.4/0.35, 0.8/0.1, -0.5/-0.3, 0.1/-0.4, 0.6/-0.35, -0.9/-0.1}
  \node[circle, fill=popblue, inner sep=1.6pt] at (\x,\y) {};
\foreach \x/\y in {-0.3/0.1, 0.3/-0.05}
  \node[circle, fill=poporange, inner sep=1.6pt] at (\x,\y) {};
\node[font=\scriptsize, text=popdark, align=center] at (0,-1.35) {$p = 0.8$ \quad $q = 0.2$};
% arrow 1
\draw[-{Stealth}, very thick, popmuted] (1.8,0) -- (3.0,0);
% Generation 2
\draw[fill=popblueL, draw=popblue, thick] (4.6,0) ellipse (1.5 and 1.0);
\node[font=\bfseries\small, text=popdark] at (4.6,1.35) {Generation 2};
\foreach \x/\y in {-0.7/0.3, -0.2/0.5, 0.4/0.35, 0.8/0.1, -0.5/-0.3, 0.1/-0.4}
  \node[circle, fill=popblue, inner sep=1.6pt] at (4.6+\x,\y) {};
\foreach \x/\y in {0.6/-0.35, -0.9/-0.1, -0.3/0.1, 0.3/-0.05}
  \node[circle, fill=poporange, inner sep=1.6pt] at (4.6+\x,\y) {};
\node[font=\scriptsize, text=popdark, align=center] at (4.6,-1.35) {$p = 0.6$ \quad $q = 0.4$};
% arrow 2
\draw[-{Stealth}, very thick, popmuted] (6.4,0) -- (7.6,0);
% Generation 3
\draw[fill=popblueL, draw=popblue, thick] (9.2,0) ellipse (1.5 and 1.0);
\node[font=\bfseries\small, text=popdark] at (9.2,1.35) {Generation 3};
\foreach \x/\y in {-0.7/0.3, -0.2/0.5, 0.4/0.35, 0.8/0.1}
  \node[circle, fill=popblue, inner sep=1.6pt] at (9.2+\x,\y) {};
\foreach \x/\y in {-0.5/-0.3, 0.1/-0.4, 0.6/-0.35, -0.9/-0.1, -0.3/0.1, 0.3/-0.05}
  \node[circle, fill=poporange, inner sep=1.6pt] at (9.2+\x,\y) {};
\node[font=\scriptsize, text=popdark, align=center] at (9.2,-1.35) {$p = 0.4$ \quad $q = 0.6$};
% legend
\node[circle, fill=popblue, inner sep=1.8pt, label={[font=\scriptsize, text=popdark]right:allele $A$}] at (2.6,-2.1) {};
\node[circle, fill=poporange, inner sep=1.8pt, label={[font=\scriptsize, text=popdark]right:allele $a$}] at (5.4,-2.1) {};
\end{tikzpicture}
\legende{A gene pool changing over three generations. Each dot represents one allele copy; the frequencies $p$ (allele $A$) and $q$ (allele $a$) shift with time. \emph{Evolution is precisely this change in allele frequencies.}}
\end{popfigure}

\begin{aretenir}[The modern definition of evolution]
\cle{Evolution}, in the modern genetic sense, is \textbf{any change in the allele frequencies of a population's gene pool over successive generations}. Individuals do not evolve; populations do.
\end{aretenir}

\courssub{Sources of Genetic Variation}

Natural selection can only act if individuals differ. Where does this \cle{genetic variation} come from? There are three great sources.

\textbf{1. Mutation --- the ultimate source of new alleles.}\par
A \cle{mutation} is a random, heritable change in the DNA sequence. Mutations are the \emph{only} source of \textbf{completely new alleles}. Most mutations are neutral or harmful, but occasionally one is advantageous --- for example, a mutation allowing a malaria parasite to break down a drug. Mutation rates are low (typically $10^{-5}$ to $10^{-6}$ per gene per generation), but over thousands of generations and millions of individuals, mutation supplies abundant raw material.

\textbf{2. Meiosis --- reshuffling existing alleles.}\par
Meiosis does not create new alleles, but it produces new \emph{combinations} of alleles in two ways:
\begin{itemize}
\item \cle{Crossing over} (prophase I): homologous chromosomes exchange segments, producing recombinant chromatids that carry new allele combinations on the same chromosome.
\item \cle{Independent assortment} (metaphase I): the orientation of each pair of homologous chromosomes is random, so with $n$ chromosome pairs there are $2^{n}$ possible combinations of maternal and paternal chromosomes in the gametes (in humans, $2^{23} \approx 8.4 \times 10^{6}$).
\end{itemize}

\textbf{3. Random fertilisation.}\par
Any sperm can fuse with any egg. This multiplies the number of possible zygotes enormously: in humans, two parents can theoretically produce about $2^{23} \times 2^{23} \approx 7 \times 10^{13}$ genetically different zygotes, even before crossing over is considered.

\begin{popfigure}
\begin{tikzpicture}[
  box/.style={draw=popblue, thick, fill=popblueL, rounded corners, align=center, font=\small, text width=3.1cm, minimum height=1.1cm},
  src/.style={draw=popgreen, thick, fill=popgreenL, rounded corners, align=center, font=\small, text width=3.4cm, minimum height=1.0cm},
  arr/.style={-{Stealth}, very thick, popmuted}]
\node[box] (var) at (0,0) {\textbf{Genetic variation} in the population};
\node[src, align=center] (mut) at (-5.2,2.2){\textbf{Mutation}\\ new alleles (ultimate source)};
\node[src, align=center] (co) at (-5.2,0){\textbf{Crossing over}\\ (prophase I of meiosis)};
\node[src, align=center] (ia) at (5.2,2.2){\textbf{Independent assortment}\\ (metaphase I)};
\node[src, align=center] (rf) at (5.2,0){\textbf{Random fertilisation}\\ any sperm + any egg};
\draw[arr] (mut.east) -- (var.165);
\draw[arr] (co.east) -- (var.195);
\draw[arr] (ia.west) -- (var.15);
\draw[arr] (rf.west) -- (var.-15);
\node[draw=poporange, thick, fill=white, rounded corners, align=center, font=\small, text width=4.6cm] (sel) at (0,-2.3) {Raw material on which \textbf{natural selection} acts};
\draw[arr, poporange] (var.south) -- (sel.north);
\end{tikzpicture}
\legende{Flow chart summarising the sources of genetic variation. Mutation creates new alleles; meiosis and random fertilisation reshuffle them into new combinations.}
\end{popfigure}

\courssub{Continuous and Discontinuous Variation}

Variation within a population falls into two broad patterns, which you must be able to distinguish with examples.

\begin{definition}[Continuous and discontinuous variation]
\cle{Discontinuous variation}: individuals fall into \textbf{distinct, separate categories} with no intermediates. It is usually controlled by \textbf{one or few genes} (monogenic or oligogenic) and is little affected by the environment. Examples: human ABO blood groups, ability to roll the tongue, albinism, sickle cell genotype.

\cle{Continuous variation}: the character shows a \textbf{complete range of intermediates} between two extremes, giving a bell-shaped (normal) distribution. It is controlled by \textbf{many genes} (\cle{polygenic} inheritance) and is strongly influenced by the \textbf{environment}. Examples: human height, skin colour, yield of maize cobs.
\end{definition}

\begin{center}
\begin{tabularx}{\linewidth}{|l|X|X|}
\hline
\rowcolor{popblueL} \textbf{Feature} & \textbf{Discontinuous variation} & \textbf{Continuous variation} \\
\hline
Categories & Distinct classes, no intermediates & Full range of intermediates \\
\hline
Genetic control & One or few genes & Many genes (polygenic) \\
\hline
Environmental effect & Little or none & Large \\
\hline
Graph shape & Separate bars & Bell-shaped curve \\
\hline
Examples & ABO blood group, sickle cell trait & Height, mass, milk yield \\
\hline
\end{tabularx}
\end{center}

\begin{attention}[Do not confuse the two types]
A character such as ``resistant vs non-resistant'' to a pesticide is discontinuous at a given moment, but the \emph{underlying} resistance level may be polygenic. In GCE answers, always justify your classification by mentioning the number of genes involved and the effect of the environment.
\end{attention}

\courssub{The Hardy-Weinberg Principle}

\textbf{The question.}\par
If nothing disturbs a population, do its allele frequencies change from generation to generation? Intuition might suggest that dominant alleles would automatically ``take over''. In 1908, Hardy and Weinberg independently showed that this is \textbf{wrong}: under ideal conditions, allele and genotype frequencies remain \textbf{constant}.

\begin{propriete}[The Hardy-Weinberg equilibrium]
In a large, randomly mating population with no mutation, no migration and no natural selection, allele and genotype frequencies remain constant from generation to generation. For a gene with two alleles $A$ and $a$ of frequencies $p$ and $q$:
\[
p + q = 1
\]
and the genotype frequencies at equilibrium are given by the binomial expansion:
\[
p^{2} + 2pq + q^{2} = 1
\]
where $p^{2} = $ frequency of $AA$, $2pq = $ frequency of $Aa$, $q^{2} = $ frequency of $aa$.

\textbf{The five conditions} (learn them by heart for the GCE):
\begin{enumerate}
\item Very large population (no genetic drift);
\item Random mating (no mate choice based on genotype);
\item No mutation;
\item No migration (no gene flow in or out);
\item No natural selection (all genotypes survive and reproduce equally).
\end{enumerate}
The derivation (a Punnett square of random gamete fusion) is instructive and may be asked: gametes carry $A$ with probability $p$ and $a$ with probability $q$, so zygotes are $AA$ with probability $p \times p = p^{2}$, $Aa$ with probability $2pq$, and $aa$ with probability $q^{2}$.
\end{propriete}

\begin{attention}[The classic trap: $p$ versus $p^{2}$]
The most common GCE mistake is to confuse the \textbf{allele frequency} $p$ (frequency of allele $A$) with the \textbf{genotype frequency} $p^{2}$ (frequency of homozygous dominant individuals $AA$). Remember: alleles come in single copies ($p$, $q$); genotypes come in pairs ($p^{2}$, $2pq$, $q^{2}$). Similarly, $q^{2}$ is the frequency of the recessive \emph{phenotype} (genotype $aa$), not of the allele $a$.
\end{attention}

\begin{methode}[Solving any Hardy-Weinberg problem]
\begin{enumerate}
\item \textbf{Start from the recessive phenotype.} Its frequency is $q^{2}$ (it is the only phenotype whose genotype is known with certainty).
\item Take the square root: $q = \sqrt{q^{2}}$.
\item Use $p = 1 - q$.
\item Compute the genotype frequencies: $p^{2}$ ($AA$), $2pq$ ($Aa$), $q^{2}$ ($aa$).
\item If asked for numbers of individuals, multiply each frequency by the population size.
\item Check: $p^{2} + 2pq + q^{2}$ must equal 1.
\end{enumerate}
\end{methode}

\begin{exoresolu}[Sickle cell frequencies in a Cameroonian population]
In a town in the Centre Region of Cameroon, a survey of 10 000 newborns finds that 400 have sickle cell anaemia (genotype $Hb^{S}Hb^{S}$). Assuming Hardy-Weinberg equilibrium, calculate (a) the frequency of the $Hb^{S}$ allele, (b) the frequency of the $Hb^{A}$ allele, (c) the percentage of the population that are carriers (heterozygotes, $Hb^{A}Hb^{S}$).
\tcblower
\textbf{Step 1 --- recessive phenotype.} Sickle cell anaemia is the homozygous recessive condition:
\[
q^{2} = \frac{400}{10\,000} = 0.04
\]
\textbf{Step 2 --- allele frequency $q$.}
\[
q = \sqrt{0.04} = 0.2
\]
\textbf{Step 3 --- allele frequency $p$.}
\[
p = 1 - q = 1 - 0.2 = 0.8
\]
\textbf{Step 4 --- carrier frequency.}
\[
2pq = 2 \times 0.8 \times 0.2 = 0.32
\]
So \textbf{32 \%} of the population are healthy carriers of the sickle cell trait (about 3 200 newborns out of 10 000), while $p^{2} = 0.64$ (64 \%) are homozygous normal. Check: $0.64 + 0.32 + 0.04 = 1$. \hfill $\blacksquare$
\end{exoresolu}

\begin{savaistu}[Why is the sickle cell allele so common in Cameroon?]
The $Hb^{S}$ allele is far more frequent in malaria-endemic regions (such as most of Cameroon) than in regions without malaria. Heterozygotes ($Hb^{A}Hb^{S}$) are partially protected against severe malaria, so natural selection \emph{maintains} the allele in the population despite its harmful effect in homozygotes. This is called \cle{balanced polymorphism} (heterozygote advantage) --- a beautiful example of natural selection acting on allele frequencies, which you will meet again in Section 2.
\end{savaistu}

\begin{popfigure}
\begin{center}
\begin{tikzpicture}
\begin{axis}[
    ybar, bar width=0.9cm,
    width=11cm, height=6.5cm,
    ymin=0, ymax=0.30,
    ylabel={Frequency of $Hb^{S}$ allele ($q$)},
    symbolic x coords={Endemic region, Non-endemic region},
    xtick=data,
    nodes near coords,
    every node near coord/.append style={font=\small\bfseries, text=popdark},
    ylabel style={font=\small},
    x tick label style={font=\small},
    axis line style={popline},
    tick style={popline},
]
\addplot[fill=poporange, draw=popdark] coordinates {(Endemic region,0.20) (Non-endemic region,0.02)};
\end{axis}
\end{tikzpicture}
\end{center}
\legende{Illustrative comparison of the frequency of the sickle cell allele $Hb^{S}$ in a malaria-endemic region and a non-endemic region. Selection by malaria maintains the allele where the disease is present (values are illustrative, not survey data).}
\end{popfigure}

\courssub{Factors that Disturb Hardy-Weinberg Equilibrium}

The Hardy-Weinberg principle describes an \emph{ideal} population. Real populations violate one or more of the five conditions, and it is precisely these violations that cause \textbf{evolution}. You must know the five disturbing factors:

\begin{enumerate}
\item \cle{Mutation}: introduces new alleles, slowly changing $p$ and $q$. Alone it is a weak force, but it supplies the raw material for all other changes.
\item \cle{Migration (gene flow)}: when individuals enter or leave a population and breed, they carry alleles with them. Gene flow tends to make neighbouring populations genetically more similar.
\item \cle{Genetic drift}: in \textbf{small} populations, chance events (which individuals happen to mate, which gametes happen to fuse) can cause large random changes in allele frequencies from one generation to the next. An allele may even be lost completely, or become \emph{fixed} (frequency 1).
\item \cle{Non-random mating}: if individuals choose mates by phenotype (e.g.\ assortative mating, or inbreeding between relatives), genotype frequencies change --- homozygotes increase, heterozygotes decrease --- even though allele frequencies may stay the same.
\item \cle{Natural selection}: genotypes differ in survival and reproductive success, so alleles that confer advantage increase in frequency. This is the only factor that produces \textbf{adaptive} change, and it is the subject of Section 2.
\end{enumerate}

\textbf{Genetic drift and the founder effect.}\par
Genetic drift is strongest in small populations. A striking case is the \cle{founder effect}: when a \textbf{small group} of individuals leaves a population and colonises a new, isolated area (an island, a remote highland valley), the new population's gene pool contains only the alleles the founders happened to carry. Rare alleles may become common, and common alleles may be absent, purely by chance. Small, isolated human communities --- for example isolated villages in the Adamawa highlands --- may therefore show unusual frequencies of certain alleles compared with the general population. The founder effect explains why some rare genetic disorders are unusually frequent in particular isolated communities worldwide.

\begin{aretenir}[Key points of Section 1]
\begin{itemize}
\item The \textbf{population} is the unit of evolution; evolution = change in \textbf{allele frequencies} of the gene pool.
\item Variation arises from \textbf{mutation} (new alleles), \textbf{crossing over} and \textbf{independent assortment} in meiosis, and \textbf{random fertilisation} (new combinations).
\item \textbf{Discontinuous} variation: few genes, distinct classes, little environmental effect. \textbf{Continuous} variation: polygenic, bell curve, strong environmental effect.
\item Hardy-Weinberg: $p + q = 1$ and $p^{2} + 2pq + q^{2} = 1$, under five strict conditions.
\item Always start calculations from $q^{2}$, the recessive phenotype frequency.
\item Mutation, migration, genetic drift, non-random mating and natural selection disturb equilibrium and drive evolution; drift and the founder effect dominate in small isolated populations.
\end{itemize}
\end{aretenir}

\begin{exercice}[GCE-style practice]
In a population of 2 000 people in the Far North Region, 18 individuals suffer from albinism, a recessive condition. Assuming Hardy-Weinberg equilibrium, calculate: (a) the frequency of the albino allele; (b) the frequency of the normal allele; (c) the expected number of carriers in the population. (d) State two conditions that this population must satisfy for your calculation to be valid.
\end{exercice}

\begin{corrige}[Exercise 1]
(a) $q^{2} = \dfrac{18}{2000} = 0.009$, so $q = \sqrt{0.009} \approx 0.095$.

(b) $p = 1 - 0.095 = 0.905$.

(c) Carrier frequency $2pq = 2 \times 0.905 \times 0.095 \approx 0.172$; number of carriers $\approx 0.172 \times 2000 \approx 344$ people.

(d) Any two of: large population, random mating, no mutation, no migration, no natural selection acting on this gene.
\end{corrige}

\coursec{Section 2: Natural Selection — Mechanism and Types}

In Section 1 we established that populations are not uniform: they contain \cle{genetic variation}, and this variation is reshuffled by meiosis and fertilisation and occasionally enriched by mutation. We also saw that the genetic make-up of a population can be described by \cle{allele frequencies}. The natural question is now: \emph{what changes allele frequencies over generations, and how?} The answer, first given a coherent form by Charles Darwin and Alfred Russel Wallace in the nineteenth century, is \cle{natural selection}. This section is the heart of the chapter: the GCE examiner expects you to explain the mechanism step by step, to recognise its three modes on a graph, and to apply it to classic examples such as antibiotic resistance and sickle cell anaemia.

\courssub{Darwin's Observations and Deductions}

Imagine a field of maize near Bafoussam. A farmer plants hundreds of seeds; only a fraction of the seedlings will survive the dry spell, the striga weed and the weevils, and of those survivors only some will produce cobs that contribute seed to the next planting. Nothing here is mysterious, yet from such everyday observations Darwin built one of the most powerful arguments in science.

\textbf{Darwin's observations.}
\begin{itemize}
\item \textbf{Observation 1 — Overproduction.} All organisms produce far more offspring than can possibly survive. A single female mosquito lays hundreds of eggs; a tilapia produces thousands.
\item \textbf{Observation 2 — Constancy of population size.} Despite this overproduction, natural populations remain roughly stable in size over long periods.
\item \textbf{Observation 3 — Variation.} Individuals within a population differ from one another, and much of this variation is heritable (Section 1).
\end{itemize}

\textbf{Darwin's deductions.}
\begin{itemize}
\item \textbf{Deduction 1 — Struggle for existence.} From Observations 1 and 2: since more offspring are produced than can survive, there is a \cle{struggle for existence} (competition for food, space, mates; exposure to predators, parasites and drought). This includes both \emph{intraspecific} competition (between members of the same species) and \emph{interspecific} competition (between different species).
\item \textbf{Deduction 2 — Survival of the fittest.} From Observation 3 and Deduction 1: individuals whose heritable traits best suit them to the environment tend to survive and reproduce more than others. This is \emph{differential survival and reproduction}. The phrase ``survival of the fittest'' was coined by Herbert Spencer but adopted by Darwin in later editions of \emph{On the Origin of Species}.
\item \textbf{Deduction 3 — Evolutionary change.} Over many generations, favourable heritable traits (and the alleles behind them) become more common in the population: the population \emph{evolves}.
\end{itemize}

\begin{definition}[Natural selection]
\cle{Natural selection} is the process by which individuals with heritable traits better suited to their environment survive and reproduce more successfully than others, so that the alleles underlying those traits increase in frequency in the population from one generation to the next.
\end{definition}

\begin{attention}[Common mistake: individuals do not evolve]
A very frequent GCE error is to write that ``the organism adapts during its lifetime'' or that ``the bacterium decided to become resistant''. This is wrong on two counts. First, the variation arises by \emph{random} mutation \emph{before} the selection pressure acts — the environment does not direct the mutation. Second, \textbf{populations evolve; individuals do not}. An individual bacterium either already carries a resistance allele or it does not; what changes over generations is the \emph{proportion} of resistant individuals in the population. Always phrase your answer in terms of changing allele frequencies.
\end{attention}

\courssub{Fitness and Selection Pressure}

\begin{definition}[Fitness and selection pressure]
In biology, \cle{fitness} means \textbf{relative reproductive success}: the contribution an individual makes to the gene pool of the next generation, relative to other individuals in the same population. It has nothing to do with physical strength. A small, weak-looking lizard that leaves twelve offspring is fitter than a muscular one that leaves none.\par
A \cle{selection pressure} is any environmental factor that reduces the survival or reproductive success of some phenotypes more than others, thereby driving natural selection.
\end{definition}

Examples of selection pressures you should be able to quote:
\begin{itemize}
\item \textbf{Predation} — camouflaged prey survive better (peppered moth, below).
\item \textbf{Disease} — \emph{Plasmodium falciparum} malaria selects for the sickle cell heterozygote in Cameroon.
\item \textbf{Antibiotics} — drugs such as amoxicillin kill susceptible bacteria, selecting resistant mutants.
\item \textbf{Pesticides} — DDT and pyrethroids select resistant \emph{Anopheles} mosquitoes.
\item \textbf{Climate} — drought selects deep-rooted or early-flowering plants.
\item \textbf{Competition for mates} — sexual selection (a subset of natural selection) favours traits that increase mating success.
\end{itemize}

\begin{popfigure}
\begin{tikzpicture}[font=\small, node distance=0.9cm]
\node[draw=popblue, fill=popblueL, rounded corners, text width=3.1cm, align=center] (var) {\textbf{Variation}\\ heritable differences in the population};
\node[draw=poporange, fill=white, rounded corners, text width=3.1cm, align=center, right=of var] (press) {\textbf{Selection pressure}\\ predation, disease, drought, antibiotics};
\node[draw=poppurple, fill=white, rounded corners, text width=3.1cm, align=center, right=of press] (diff) {\textbf{Differential survival}\\ fittest phenotypes reproduce more};
\node[draw=popgreen, fill=popgreenL, rounded corners, text width=3.4cm, align=center, right=of diff] (freq) {\textbf{Allele frequency change}\\ population evolves over generations};
\draw[-latex, thick, popdark] (var) -- (press);
\draw[-latex, thick, popdark] (press) -- (diff);
\draw[-latex, thick, popdark] (diff) -- (freq);
\end{tikzpicture}
\legende{The logical chain of natural selection: heritable variation, acted on by a selection pressure, produces differential survival and reproduction, which changes allele frequencies.}
\end{popfigure}

\courssub{The Three Types of Natural Selection}

Natural selection can push a population's trait distribution in three different ways. You must be able to \emph{draw} and \emph{interpret} all three.

\textbf{1. Directional selection.} One extreme phenotype is favoured, so the mean of the population shifts in one direction over generations. Classic examples: \cle{antibiotic resistance} in bacteria (the most resistant cells survive treatment and multiply), \cle{pesticide resistance} in \emph{Anopheles} mosquitoes, and the peppered moth (below).

\textbf{2. Stabilising selection.} The intermediate phenotype is favoured and both extremes are eliminated, so the distribution becomes narrower. Classic example: \cle{human birth weight}. Babies that are too small are vulnerable; babies that are too large risk complications at birth. Babies of intermediate weight (around $3$--$3.5\ \mathrm{kg}$) survive best. Another example is \cle{clutch size} in birds: laying too few eggs wastes reproductive potential, laying too many means none of the chicks can be fed adequately.

\textbf{3. Disruptive selection.} Both extremes are favoured at the expense of the intermediate phenotype, so the distribution may split into two peaks. Classic example: \cle{beak size} in seed-eating birds (such as the African finch \emph{Pyrenestes}) where only small soft seeds and large hard seeds are available — small-beaked birds handle small seeds, large-beaked birds crack large seeds, but medium beaks are good at neither.

\begin{popfigure}
\begin{tikzpicture}
\begin{axis}[width=15.5cm, height=5.2cm, axis lines=left, xlabel={Phenotype}, ylabel={Frequency}, xtick=\empty, ytick=\empty, legend style={font=\scriptsize, at={(1.02,1)}, anchor=north west}, domain=0:10, samples=80]
\addplot[popmuted, dashed, thick] {exp(-((x-5)^2)/2)}; \addlegendentry{Before selection}
\addplot[poporange, very thick] {exp(-((x-7)^2)/2)}; \addlegendentry{Directional}
\addplot[popblue, very thick] {1.6*exp(-((x-5)^2)/0.6)}; \addlegendentry{Stabilising}
\addplot[popgreen, very thick] {0.8*exp(-((x-2.6)^2)/1.2)+0.8*exp(-((x-7.4)^2)/1.2)}; \addlegendentry{Disruptive}
\end{axis}
\end{tikzpicture}
\legende{The three modes of selection on an originally bell-shaped (normal) distribution. Directional: the mean shifts. Stabilising: the curve narrows around the mean. Disruptive: two peaks emerge at the extremes.}
\end{popfigure}

\begin{methode}[Interpreting a selection graph]
\begin{enumerate}
\item Identify the trait on the $x$-axis (e.g. birth weight, beak size) and frequency on the $y$-axis.
\item Look at the ``before'' curve and the ``after'' curve (or the shaded survival zone).
\item Ask: \emph{which phenotype survives best?} One extreme $\rightarrow$ directional; the intermediate $\rightarrow$ stabilising; both extremes $\rightarrow$ disruptive.
\item State the consequence for the population: shift of the mean, narrowing of the curve, or splitting into two peaks.
\item Conclude in terms of \textbf{allele frequencies}: alleles for the favoured phenotype increase over generations.
\end{enumerate}
\end{methode}

\courssub{Worked Example: The Peppered Moth}

\begin{experience}[Case study: \emph{Biston betularia}, a directional selection classic]
Before the industrial revolution in England, the light, speckled form of the peppered moth was camouflaged on lichen-covered tree bark, while the rare dark (melanic) mutant was easily spotted and eaten by birds. During industrialisation, soot killed the lichens and blackened the bark: the dark form became camouflaged and the light form conspicuous, so the dark allele rose to high frequency near industrial cities. Bernard Kettlewell's experiments in the 1950s confirmed bird predation as the selective agent. After clean-air laws reduced pollution in the twentieth century, the trend reversed. The selection pressure throughout was \textbf{predation by birds}; the change was a shift in the frequency of the melanic allele — a beautiful, directly observed case of directional selection.
\end{experience}

\begin{popfigure}
\begin{tikzpicture}[font=\scriptsize]
% Tree trunks
\fill[popblueL] (0,0) rectangle (3.2,2.6);
\fill[popmuted!60] (4.2,0) rectangle (7.4,2.6);
\fill[popblueL] (8.4,0) rectangle (11.6,2.6);
% Moths: light = white ellipse, dark = black ellipse
\fill[white, draw=popdark] (1.0,1.8) ellipse (0.28 and 0.16);
\fill[white, draw=popdark] (2.2,0.9) ellipse (0.28 and 0.16);
\fill[popdark] (1.6,1.2) ellipse (0.28 and 0.16);
\fill[popdark] (4.9,1.7) ellipse (0.28 and 0.16);
\fill[popdark] (6.3,0.8) ellipse (0.28 and 0.16);
\fill[white, draw=popdark] (5.6,1.1) ellipse (0.28 and 0.16);
\fill[white, draw=popdark] (9.2,1.7) ellipse (0.28 and 0.16);
\fill[white, draw=popdark] (10.6,0.9) ellipse (0.28 and 0.16);
\fill[popdark] (9.9,1.2) ellipse (0.28 and 0.16);
\node[align=center] at (1.6,-0.5) {Before pollution\\ light form common};
\node[align=center] at (5.8,-0.5) {Industrial era\\ dark form common};
\node[align=center] at (10.0,-0.5) {Clean air restored\\ light form returns};
\end{tikzpicture}
\legende{Directional selection in the peppered moth: bird predation favours whichever morph is camouflaged, so the melanic allele frequency rises with pollution and falls when pollution is removed.}
\end{popfigure}

\begin{exoresolu}[Antibiotic resistance in bacteria]
\textbf{Statement.} A patient in Yaoundé is treated with an antibiotic for a bacterial infection. After a few days most bacteria are killed, but the infection later returns and no longer responds to the same antibiotic. Explain, using the theory of natural selection, how resistance arose. (4 marks)
\tcblower
\textbf{Solution.}
\begin{enumerate}
\item The bacterial population contains \textbf{variation}: by random mutation, a few bacteria carry an allele conferring resistance to the antibiotic. The mutation occurs \emph{spontaneously}, not because of exposure to the antibiotic. (1 mark)
\item The antibiotic acts as a \textbf{selection pressure}: susceptible bacteria are killed, while resistant bacteria survive. (1 mark)
\item The surviving resistant bacteria \textbf{reproduce} and pass the resistance allele to their offspring (differential survival and reproduction). (1 mark)
\item Over generations the \textbf{frequency of the resistance allele increases} in the population, so the population becomes resistant — the population has evolved. (1 mark)
\end{enumerate}
Note the forbidden phrasing: ``the bacteria became resistant in order to survive'' would score nothing, because it suggests directed adaptation (Lamarckism).
\end{exoresolu}

\courssub{Artificial Selection versus Natural Selection}

Humans have practised selection for millennia. In \cle{artificial selection} (also called \cle{selective breeding}), \emph{humans} choose which individuals reproduce, selecting for traits useful or pleasing to us. Cameroonian farmers practise it whenever they keep seed from the highest-yielding maize plants, or select cassava cuttings from plants resistant to cassava mosaic disease, or breed from the bull that produces the most meat or the cow giving the most milk. The result over generations is improved crop varieties and livestock breeds.

\begin{center}
\begin{tabularx}{\linewidth}{|l|X|X|}
\hline
\rowcolor{popblueL} & \textbf{Natural selection} & \textbf{Artificial selection} \\
\hline
Agent of selection & The environment (selection pressures) & Humans \\
\hline
Trait favoured & Survival and reproductive success in nature & Traits useful or desirable to humans \\
\hline
Speed & Usually slow (many generations) & Can be rapid under strict breeding \\
\hline
Example & Peppered moth, antibiotic resistance & Improved maize, dairy cattle, cassava varieties \\
\hline
\end{tabularx}
\end{center}

The underlying mechanism is identical — differential reproduction changing allele frequencies — only the ``selector'' differs.

\courssub{Balancing Selection: Sickle Cell Anaemia and Malaria}

Sometimes selection does not push a population towards a single best genotype but \emph{maintains two alleles} in the population. This is \cle{balancing selection}, and its best-documented form is \cle{heterozygote advantage} (also called \cle{heterosis}).

\begin{experience}[Case study: sickle cell heterozygote advantage in Cameroon]
The allele $H\!b^{S}$ for sickle cell haemoglobin is common in the forest and savanna zones of Cameroon, where \emph{Plasmodium falciparum} malaria is endemic. Consider the three genotypes:
\begin{itemize}
\item $H\!b^{A}H\!b^{A}$: normal haemoglobin, but fully susceptible to severe malaria — reduced fitness in malaria zones.
\item $H\!b^{A}H\!b^{S}$ (heterozygote, sickle cell \emph{trait}): phenotypically almost normal, and the parasite develops poorly in their red blood cells — \textbf{highest fitness} where malaria is present.
\item $H\!b^{S}H\!b^{S}$: sickle cell \emph{anaemia}, a serious disease often fatal without treatment — very low fitness.
\end{itemize}
Because the heterozygote is fitter than \emph{both} homozygotes, neither allele is eliminated: selection keeps $H\!b^{S}$ at a relatively high frequency precisely where malaria is common. This explains the geographical overlap between the sickle cell allele and malaria across sub-Saharan Africa — a classic example of a \cle{balanced polymorphism}.
\end{experience}

\begin{exoresolu}[Calculating genotype frequencies]
\textbf{Statement.} In a population in a malaria zone, the frequency of the $H\!b^{S}$ allele is $q = 0.2$. Assuming Hardy--Weinberg equilibrium, calculate the frequencies of the three genotypes and comment.
\tcblower
\textbf{Solution.} With $p = 1 - q = 0.8$:
\begin{align*}
H\!b^{A}H\!b^{A} &: p^{2} = 0.8^{2} = 0.64 \\
H\!b^{A}H\!b^{S} &: 2pq = 2 \times 0.8 \times 0.2 = 0.32 \\
H\!b^{S}H\!b^{S} &: q^{2} = 0.2^{2} = 0.04
\end{align*}
Comment: nearly one person in three is a protected heterozygote, while about $4\%$ suffer sickle cell anaemia. The high frequency of a harmful allele is maintained only because heterozygote advantage (balancing selection) outweighs the loss of $H\!b^{S}H\!b^{S}$ individuals.
\end{exoresolu}

\begin{attention}[Exam tip: always finish at allele frequencies]
Whatever the example — moths, bacteria, mosquitoes, birth weight — the examiner's mark scheme always contains a final point equivalent to: ``the frequency of the favourable allele increases over generations''. Structure every natural-selection answer as: variation $\rightarrow$ selection pressure $\rightarrow$ differential survival/reproduction $\rightarrow$ \textbf{change in allele frequency}. Also remember: fitness = reproductive success, never physical strength.
\end{attention}

\begin{exercice}[Applying the three types of selection]
\begin{enumerate}
\item For each situation, name the type of selection and justify: (a) in a drought-prone region, only the deepest-rooted seedlings survive; (b) human babies of intermediate birth weight have the lowest infant mortality; (c) in a habitat containing only very small and very large seeds, finches with medium beaks starve.
\item Explain why pyrethroid-treated bed nets, used alone for many years, may become less effective against \emph{Anopheles} mosquitoes.
\item State two differences and one similarity between natural and artificial selection.
\end{enumerate}
\end{exercice}

\begin{corrige}[Exercise 1]
\begin{enumerate}
\item (a) \textbf{Directional}: one extreme (deep roots) is favoured, the mean shifts. (b) \textbf{Stabilising}: the intermediate phenotype is favoured, extremes eliminated. (c) \textbf{Disruptive}: both extremes favoured, intermediates eliminated.
\item The net's insecticide is a selection pressure. Rare mosquitoes carrying resistance alleles survive exposure and reproduce; susceptible mosquitoes die. Over generations the frequency of resistance alleles rises, so the mosquito population becomes resistant and the nets lose effectiveness.
\item Differences: the agent (environment vs humans) and the goal (survival/reproduction vs human usefulness). Similarity: both act through differential reproduction causing allele frequencies to change over generations.
\end{enumerate}
\end{corrige}

\begin{aretenir}[Key points of Section 2]
\begin{itemize}
\item Natural selection: overproduction + heritable variation $\rightarrow$ struggle for existence $\rightarrow$ differential survival and reproduction $\rightarrow$ allele frequency change.
\item \textbf{Fitness = relative reproductive success}, not strength. \textbf{Populations evolve; individuals do not.}
\item A selection pressure is any environmental factor (predation, disease, antibiotics, climate, competition for mates) favouring some phenotypes.
\item Three modes: \textbf{directional} (mean shifts), \textbf{stabilising} (curve narrows), \textbf{disruptive} (two peaks).
\item Artificial selection uses the same mechanism with humans as the selector.
\item Sickle cell trait in malaria zones is \textbf{balancing selection by heterozygote advantage} (balanced polymorphism).
\end{itemize}
\end{aretenir}

\coursec{Section 3: Speciation and Isolating Mechanisms}

In Section 2 we saw how \cle{natural selection} acts on variation within a population, changing allele frequencies from generation to generation. But natural selection alone only explains how a population \emph{changes}; it does not yet explain how one species gives rise to \emph{two} species. If you travel from the lowland rainforest around Kumba up the slopes of Mount Cameroon, you will notice that related plants and animals at different altitudes can no longer interbreed. How does a single interbreeding population split into groups that are reproductively separate? That is the question of \cle{speciation}, and it is the bridge between microevolution (changes within populations) and macroevolution (the origin of new species).

\courssub{The Biological Species Concept}

\begin{definition}[Species (biological species concept)]
A \cle{species} is a group of actually or potentially interbreeding natural populations whose members can mate with one another and produce \textbf{fertile offspring}, and which are \textbf{reproductively isolated} from other such groups.
\end{definition}

The two key tests are: (1) can the individuals \emph{interbreed} under natural conditions, and (2) are the offspring \emph{fertile}? A horse and a donkey can mate and produce a mule, but the mule is sterile, so horse and donkey remain \emph{separate species}. Two goats, one from a farm in Bamenda and one from Buea, may look different, yet they interbreed freely and produce fertile kids: they belong to the \emph{same} species.

\begin{attention}[Limits of the biological species concept]
The concept works well for sexually reproducing organisms but fails for asexual organisms (many bacteria, some plants) and for fossils, where interbreeding cannot be tested. In such cases morphological or genetic criteria are used. In GCE essays, always state the biological species concept as the standard definition, then mention its limits only if asked.
\end{attention}

\begin{definition}[Speciation]
\cle{Speciation} is the evolutionary process by which one species splits into two or more new species. It occurs when populations of an ancestral species become \textbf{reproductively isolated}, so that gene flow between them stops and they accumulate independent genetic differences by mutation, selection and genetic drift.
\end{definition}

The central idea is \cle{reproductive isolation}. As long as genes flow between two populations, any advantageous allele can spread through both, and the populations evolve together as one species. Once gene flow is cut, each population evolves independently; given enough time, differences accumulate until even if the populations meet again they can no longer produce fertile offspring. Speciation is then complete and irreversible.

\courssub{Allopatric Speciation}

\begin{definition}[Allopatric speciation]
\cle{Allopatric speciation} (from Greek \emph{allos} = other, \emph{patra} = homeland) is speciation that occurs when a population is split by a \textbf{geographical barrier} (mountain range, river, sea, desert, lava flow) into two or more isolated populations which then diverge genetically until they become separate species.
\end{definition}

The logic is intuitive: imagine a single population of forest birds spread across a valley. A river changes course, or volcanic activity raises a ridge, splitting the population in two. The two halves face different climates, different foods, different predators. Natural selection favours different alleles on each side; different mutations arise; drift acts differently. After thousands of generations, the two populations are so different that even if the barrier disappears, they can no longer interbreed successfully. Two species now exist where there was one.

\begin{popfigure}
\begin{center}
\begin{tikzpicture}[font=\small, >=stealth]
% Stage 1: single population
\draw[fill=popgreenL, draw=popgreen, thick] (0,0) ellipse (1.6 and 0.8);
\node at (0,0) {\textbf{One population}};
\node at (0,-1.2) {free gene flow};
% arrow 1
\draw[->, very thick, poporange] (2.1,0) -- (3.1,0) node[midway, above, text=popdark]{barrier forms};
% Stage 2: barrier + two populations
\draw[fill=popgreenL, draw=popgreen, thick] (4.6,0.6) ellipse (1.0 and 0.55);
\draw[fill=popblueL, draw=popblue, thick] (4.6,-0.9) ellipse (1.0 and 0.55);
\draw[fill=popmuted, draw=popdark] (5.9,-1.7) rectangle (6.3,1.4);
\node[text=popdark] at (6.1,1.7) {barrier};
\node at (4.6,0.6) {Pop.\ A};
\node at (4.6,-0.9) {Pop.\ B};
\node at (4.6,-2.1) {no gene flow};
% arrow 2
\draw[->, very thick, poporange] (6.7,0) -- (7.7,0) node[midway, above, text=popdark, text width=2.2cm, align=center]{independent evolution};
% Stage 3: two species
\draw[fill=popgreenL, draw=popgreen, thick] (9.6,0.6) ellipse (1.0 and 0.55);
\draw[fill=poppink!25, draw=poppink, thick] (9.6,-0.9) ellipse (1.0 and 0.55);
\node at (9.6,0.6) {Species 1};
\node at (9.6,-0.9) {Species 2};
\node[text=popdark, text width=3.2cm, align=center] at (9.6,-2.1) {reproductively isolated even if barrier removed};
\end{tikzpicture}
\end{center}
\legende{Allopatric speciation: a geographical barrier stops gene flow; the isolated populations diverge until reproductive isolation is complete.}
\end{popfigure}

\begin{savaistu}[The Cameroon volcanic line]
The Cameroon volcanic line is a chain of volcanoes running from the Gulf of Guinea (Bioko, Pr\'incipe, S\~ao Tom\'e) through Mount Cameroon, Mount Manengouba, the Bamboutos Mountains and Mount Oku. These mountains and the sea channels between the islands act as geographical barriers. Populations of birds, reptiles and plants isolated on different mountains or islands have diverged, giving the region its remarkable \cle{endemism} --- species found nowhere else on Earth, such as several chameleons and the Mount Cameroon francolin. This is allopatric speciation written into our own landscape.
\end{savaistu}

\courssub{Sympatric Speciation}

\begin{definition}[Sympatric speciation]
\cle{Sympatric speciation} (Greek \emph{sym} = together) is speciation occurring \textbf{within the same geographical area}, without any physical barrier, when reproductive isolation arises through genetic, ecological or behavioural mechanisms.
\end{definition}

The most important and best-documented route is \cle{polyploidy} in plants. If an error in meiosis produces diploid gametes, self-fertilisation can yield a \emph{tetraploid} ($4n$) plant in a single generation. A tetraploid crossed with the original diploid ($2n$) gives a triploid ($3n$) offspring, which is usually sterile because its chromosomes cannot pair properly at meiosis. The tetraploid is therefore \emph{instantly} reproductively isolated from its diploid ancestors, even though it grows in the same field: a new species in one step. Many crop plants are polyploids --- bread wheat ($6n$), bananas, coffee, sugar cane.

Two forms of polyploidy are distinguished:
\begin{itemize}
\item \textbf{Autopolyploidy:} chromosome doubling within a single species (e.g.\ a diploid plant becoming tetraploid).
\item \textbf{Allopolyploidy:} chromosome doubling in a \emph{hybrid} between two species. The sterile hybrid becomes fertile after doubling, because each chromosome now has a homologous partner at meiosis. This is how bread wheat and many cultivated brassicas arose.
\end{itemize}

Other sympatric routes include:
\begin{itemize}
\item \textbf{Ecological isolation:} sub-populations exploit different micro-habitats or host plants in the same area and rarely meet to mate (e.g.\ insects shifting to a new host plant).
\item \textbf{Behavioural isolation:} mutations affecting courtship song, colour or mating time cause assortative mating within one area.
\end{itemize}

\begin{methode}[Distinguishing allopatric from sympatric speciation in exam scenarios]
When a GCE question describes the origin of two species, ask these questions in order:
\begin{enumerate}
\item \textbf{Was there a physical barrier?} If a river, mountain, sea or distance separated the populations first, the answer is \emph{allopatric}.
\item \textbf{Did the populations live in the same area throughout?} If yes, the answer is \emph{sympatric}.
\item \textbf{If sympatric, what mechanism?} Look for clues: sudden chromosome doubling in a plant $\Rightarrow$ polyploidy; different hosts or flowering times $\Rightarrow$ ecological/temporal isolation; different courtship $\Rightarrow$ behavioural isolation.
\item \textbf{Always justify} your answer with one sentence citing the evidence from the scenario --- marks are awarded for the justification, not just the label.
\end{enumerate}
\end{methode}

\courssub{Isolating Mechanisms}

Reproductive isolation is maintained by \cle{isolating mechanisms}, classified according to whether they act \emph{before} or \emph{after} fertilisation (zygote formation).

\begin{popfigure}
\begin{center}
\begin{tikzpicture}[font=\small]
\node[fill=popblue, text=white, rounded corners, text width=4.2cm, align=center] (root) at (0,0) {\textbf{Reproductive isolating mechanisms}};
\node[fill=popgreenL, draw=popgreen, rounded corners, text width=5.2cm, align=left] (pre) at (-4.6,-2.2) {\textbf{Pre-zygotic} (prevent fertilisation)\\ \textbullet\ habitat isolation\\ \textbullet\ seasonal / temporal\\ \textbullet\ behavioural\\ \textbullet\ mechanical\\ \textbullet\ gametic};
\node[fill=poppink!25, draw=poppink, rounded corners, text width=5.2cm, align=left] (post) at (4.6,-2.2) {\textbf{Post-zygotic} (zygote forms but fails)\\ \textbullet\ hybrid inviability\\ \textbullet\ hybrid sterility (mule)\\ \textbullet\ hybrid breakdown};
\draw[->, thick, popdark] (root) -- (pre);
\draw[->, thick, popdark] (root) -- (post);
\end{tikzpicture}
\end{center}
\legende{Classification of reproductive isolating mechanisms.}
\end{popfigure}

\textbf{Pre-zygotic isolating mechanisms}\par
These prevent mating or fertilisation, so no hybrid zygote is ever formed. They are energetically ``cheaper'' for the organisms because no gametes are wasted.

\begin{center}
\begin{tabularx}{\linewidth}{|l|X|X|}
\hline
\rowcolor{popblueL} \textbf{Mechanism} & \textbf{How it works} & \textbf{Example} \\
\hline
Habitat (ecological) & Species live in different habitats in the same region and rarely meet & Two antelope species, one in forest, one in savanna \\
\hline
Seasonal / temporal & Breeding seasons or daily activity times differ & Two frog species breeding in different months; plants flowering at different times \\
\hline
Behavioural & Courtship rituals, songs or displays differ; mates not recognised & Different songs of related sunbird species; firefly flash patterns \\
\hline
Mechanical & Genitalia or floral structures do not fit & Insect genitalia of related species are structurally incompatible \\
\hline
Gametic & Gametes meet but cannot fuse; pollen or sperm fails & Pollen of one plant species cannot germinate on the stigma of another; sea urchin sperm cannot fertilise eggs of another species \\
\hline
\end{tabularx}
\end{center}

\textbf{Post-zygotic isolating mechanisms}\par
Here mating and fertilisation occur, but the hybrid is eliminated or cannot reproduce.

\begin{itemize}
\item \textbf{Hybrid inviability:} the hybrid zygote fails to develop or dies before reproductive age, because the two parental genomes are incompatible.
\item \textbf{Hybrid sterility:} the hybrid survives and is vigorous but cannot produce functional gametes. The classic example is the \cle{mule} (mare, $2n = 64$, $\times$ donkey, $2n = 62$): the hybrid has 63 chromosomes which cannot pair evenly at meiosis, so the mule is sterile. The same applies to the hinny and to many plant hybrids.
\item \textbf{Hybrid breakdown:} the first-generation (F$_1$) hybrids are viable and fertile, but the F$_2$ or backcross generations are weak, deformed or sterile. This occurs in some cotton and rice crosses.
\end{itemize}

\begin{attention}[Common mistake: geographical isolation is NOT reproductive isolation]
Geographical isolation is a \emph{physical} separation; reproductive isolation is a \emph{biological} inability to produce fertile offspring. A barrier only \emph{causes} speciation if it leads, over time, to reproductive isolation. Two populations separated by a mountain but still able to interbreed if brought together are \textbf{one species}. Never write ``the mountain made them two species'' --- write ``the mountain stopped gene flow, allowing the populations to diverge until reproductive isolation evolved''.
\end{attention}

\courssub{Adaptive Radiation: Darwin's Finches}

\begin{definition}[Adaptive radiation]
\cle{Adaptive radiation} is the rapid evolution of many new species from a single ancestral species, as populations colonise and adapt to different available habitats or niches, typically after a colonisation event or the disappearance of competitors.
\end{definition}

The classic model is \cle{Darwin's finches} on the Gal\'apagos Islands. From one ancestral seed-eating finch that reached the islands from mainland South America, about 14 species evolved, differing mainly in \textbf{beak size and shape}, each adapted to a different diet: large crushing beaks for hard seeds, slender probing beaks for insects, pointed beaks for cactus flowers, and even a ``tool-using'' woodpecker finch. On the islands there were no competing birds, so each population that specialised on an unexploited food source was favoured by natural selection; geographical isolation between islands then fixed the differences into separate species. Adaptive radiation therefore combines natural selection (Section 2) with the isolating mechanisms of this section.

\begin{popfigure}
\begin{center}
\begin{tikzpicture}[font=\small, >=stealth, grow'=right, level distance=3.4cm,
 level 1/.style={sibling distance=3.6cm},
 level 2/.style={sibling distance=1.7cm}]
\node[fill=popgold!40, draw=popdark, rounded corners, text width=2.6cm, align=center] {Ancestral seed-eating finch}
 child { node[fill=popgreenL, draw=popgreen, rounded corners, text width=2.9cm, align=center] {Warbler finch\\ thin pointed beak --- insects} }
 child { node[fill=popblueL, draw=popblue, rounded corners, text width=2.9cm, align=center] {Tree finches\\ grasping beaks --- insects, nectar} }
 child { node[fill=poppink!25, draw=poppink, rounded corners, text width=2.9cm, align=center] {Ground finches}
   child { node[fill=white, draw=popline, rounded corners, text width=2.7cm, align=center] {Large ground finch\\ massive beak --- hard seeds} }
   child { node[fill=white, draw=popline, rounded corners, text width=2.7cm, align=center] {Cactus finch\\ long beak --- cactus seeds} }
 };
\end{tikzpicture}
\end{center}
\legende{Adaptive radiation of Darwin's finches: one ancestor, many beak shapes, each matched to a different food niche.}
\end{popfigure}

\begin{savaistu}[Convergent evolution --- a useful extra for essays]
\cle{Convergent evolution} is the independent evolution of similar features in unrelated species facing similar ways of life. The structures produced are \cle{analogous}: same function, different origin (e.g.\ the wing of a bird and the wing of an insect; the streamlined body of a shark, a fish, and a dolphin, a mammal). Contrast with \emph{homologous} structures (same origin, possibly different functions, e.g.\ the pentadactyl limb), which indicate common ancestry. Examiners reward candidates who can make this distinction --- it also anticipates the evidence of evolution in Section 4.
\end{savaistu}

\courssub{Worked Example and Exam Practice}

\begin{exoresolu}[Analysing a speciation scenario]
\textbf{Statement.} A single species of flowering plant grows on both sides of the Sanaga River. Geological changes widen the river so that seeds can no longer cross. After many thousands of years, botanists cross plants from the two banks: seeds form, but the seedlings always die before flowering. (a) Identify the type of speciation. (b) Name and classify the isolating mechanism observed. (c) Explain the role of the river in the process.
\tcblower
\textbf{Solution.}
(a) \emph{Allopatric speciation}: the populations were first separated by a geographical barrier (the widened river), then diverged.

(b) The hybrids die before flowering, so the mechanism is \emph{hybrid inviability}, a \emph{post-zygotic} isolating mechanism (a zygote is formed but does not survive to reproduce).

(c) The river did not ``create'' the new species directly. It acted as a geographical barrier that \emph{stopped gene flow} between the two populations. Each population then accumulated independent mutations and was exposed to different selection pressures, until genetic divergence was so great that hybrids were no longer viable. At that point reproductive isolation --- and therefore speciation --- was complete.
\end{exoresolu}

\begin{exercice}[GCE-style practice]
\begin{enumerate}
\item Define \emph{species} according to the biological species concept, and state one limitation of this concept.
\item A plant breeder crosses a diploid cabbage with a related diploid species and obtains a sterile hybrid. Treatment of the hybrid with colchicine doubles its chromosome number, and the resulting plant is fertile but cannot breed with either parent. (a) Name the type of speciation. (b) Explain why the chromosome-doubled plant is a new species.
\item Classify each of the following as pre-zygotic or post-zygotic: (i) two toad species mate at different seasons; (ii) a mule is sterile; (iii) the genitalia of two beetle species do not fit; (iv) F$_1$ hybrids are fertile but their offspring are weak.
\item Explain, with a named example, what is meant by adaptive radiation.
\end{enumerate}
\end{exercice}

\begin{corrige}[Exercise 1]
\begin{enumerate}
\item A species is a group of actually or potentially interbreeding natural populations producing fertile offspring and reproductively isolated from other such groups. Limitation: it cannot be applied to asexual organisms or to fossils.
\item (a) Sympatric speciation by allopolyploidy (chromosome doubling in a hybrid). (b) The doubled plant has paired homologous chromosomes, so meiosis is normal and it is fertile; but crosses with the diploid parents give triploid, sterile offspring. It is therefore reproductively isolated from both parents --- a new species formed in one generation, in the same area.
\item (i) pre-zygotic (temporal); (ii) post-zygotic (hybrid sterility); (iii) pre-zygotic (mechanical); (iv) post-zygotic (hybrid breakdown).
\item Adaptive radiation is the rapid production of many species from one ancestor as they adapt to different niches. Example: Darwin's finches on the Gal\'apagos, where one seed-eating ancestor gave rise to species with beaks specialised for seeds, insects, cactus flowers, etc.
\end{enumerate}
\end{corrige}

\begin{aretenir}[Exam tip for GCE essays on speciation]
A top-mark essay on speciation must contain: (1) the biological definition of a species; (2) the central role of \emph{reproductive isolation} and the stopping of gene flow; (3) a clear distinction between allopatric and sympatric speciation, each with a named example (Cameroon volcanic line or Gal\'apagos for allopatric; polyploid wheat or cabbage for sympatric); (4) \textbf{at least two named isolating mechanisms with examples} (e.g.\ behavioural isolation in bird song; hybrid sterility in the mule); (5) a link to natural selection and, if relevant, adaptive radiation. Structure the essay in this order and you will cover every marking point.
\end{aretenir}

We have now seen how new species arise. But what is the \emph{evidence} that evolution and speciation have actually occurred over the history of life? That is the subject of Section 4: the evidence of organic evolution, from fossils to molecules.

\coursec{Section 4: Evidence of Organic Evolution}

In Section 3 we saw how populations, once genetically isolated, diverge until they can no longer interbreed: speciation is the visible tip of a much deeper process. But a sceptical examiner — and the GCE examiner often plays this role — will ask: \emph{how do we actually know that evolution has happened?} After all, no one watched a reptile turn into a bird. The answer is that evolution, like a court case, is established by the \cle{convergence of many independent lines of evidence}. No single proof is decisive; but when fossils, anatomy, embryos, molecules and geography all tell the same story, the conclusion becomes overwhelming. This section examines each line of evidence in turn, in the depth the GCE Advanced Level requires.

\courssub{Fossil Evidence}

\textbf{Observation first.} In 1909, workers cutting a road through sedimentary rock near Bamenda might have noticed shells embedded in stone far from the sea. How did marine shells arrive on a hillside? The answer — that the rock was once the floor of an ancient sea, and the shells are the preserved remains of organisms that died millions of years ago — is the starting point of palaeontology.

\begin{definition}[Fossil]
A \cle{fossil} is any preserved remains, impression or trace of an organism that lived in the geological past. Fossils include mineralised bones and shells, impressions and casts in sedimentary rock, organisms trapped whole in amber or ice, and trace fossils such as footprints and burrows. The study of fossils is \cle{palaeontology}.
\end{definition}

\textbf{How fossils form.} Fossilisation is a rare event requiring rapid burial and protection from decay and scavengers. The typical sequence is:

\begin{enumerate}
\item An organism dies and its body sinks into soft sediment (mud, sand, volcanic ash), usually in water.
\item The soft parts decompose or are eaten; the hard parts (bones, teeth, shells) are buried by successive layers of sediment.
\item Over immense time, the sediment is compacted into sedimentary rock; minerals dissolved in groundwater gradually replace the original material of the hard parts (\cle{petrification}).
\item Later, erosion or human excavation exposes the fossil at the surface.
\end{enumerate}

Because sedimentary rocks are laid down in layers (\cle{strata}), the deeper a stratum, the older it is. This is the principle of \cle{superposition}. Relative ages of fossils are read from their stratum; absolute ages are obtained by \cle{radiometric dating} of the surrounding rock, using the known half-lives of isotopes such as \ce{^{14}C} (for recent remains) or potassium--argon (for ancient rocks).

\textbf{The fossil record.} When fossils are arranged in chronological order, they reveal a clear pattern: the oldest rocks contain only simple organisms (bacteria, then simple invertebrates); progressively younger rocks contain fish, then amphibians, then reptiles, then birds and mammals. This ordered succession is exactly what evolution predicts and is very hard to explain otherwise.

\begin{popfigure}
\begin{center}
\begin{tikzpicture}[scale=0.95, every node/.style={font=\small}]
% timeline axis
\draw[->, thick, popink] (0,0) -- (12.5,0) node[right]{present};
% era blocks
\fill[popblueL] (0,0) rectangle (4,1.2);
\fill[popgreenL] (4,0) rectangle (8.5,1.2);
\fill[popgold!40] (8.5,0) rectangle (12.3,1.2);
\node[align=center] at (2,0.6) {Palaeozoic\\ \scriptsize 570--245 mya};
\node[align=center] at (6.25,0.6) {Mesozoic\\ \scriptsize 245--65 mya};
\node[align=center] at (10.4,0.6) {Cenozoic\\ \scriptsize 65 mya--present};
% fossil appearances
\draw[->, popdark] (1,0) -- (1,-0.7) node[below, align=center]{first\\ invertebrates};
\draw[->, popdark] (2.6,0) -- (2.6,-1.6) node[below, align=center]{first fishes};
\draw[->, popdark] (4.6,0) -- (4.6,-0.7) node[below, align=center]{amphibians};
\draw[->, popdark] (6,0) -- (6,-1.6) node[below, align=center]{reptiles\\ (dinosaurs)};
\draw[->, popdark] (7.6,0) -- (7.6,-0.7) node[below, align=center]{\emph{Archaeopteryx}\\ (first bird)};
\draw[->, popdark] (9.6,0) -- (9.6,-1.6) node[below, align=center]{mammals\\ diversify};
\draw[->, popdark] (11.8,0) -- (11.8,-0.7) node[below, align=center]{humans};
\end{tikzpicture}
\end{center}
\legende{Simplified geological time scale showing the order of appearance of major groups in the fossil record (mya = millions of years ago).}
\end{popfigure}

\textbf{Transitional forms.} Critics once claimed that if evolution were true, we should find ``missing links'' — fossils intermediate between major groups. In fact, many such \cle{transitional forms} are known:

\begin{itemize}
\item \emph{\cle{Archaeopteryx}} (about 150 million years old) is a mosaic between reptiles and birds: it had feathers and wings (bird characters) but also teeth, clawed fingers on the wings and a long bony tail (reptilian characters). It is exactly the kind of intermediate evolution predicts.
\item The \cle{horse series} shows gradual change over about 55 million years: from the small, forest-dwelling \emph{Eohippus} (four toes on the front feet, low-crowned teeth for browsing) through intermediate genera to the modern \emph{Equus} (large, single-toed, high-crowned teeth for grazing on abrasive grass). Limb, tooth and body-size changes can be followed step by step.
\item Fossil series also link fish to amphibians (lobe-finned fish with limb-like fins) and reptiles to mammals (mammal-like reptiles).
\end{itemize}

\begin{attention}[Limits of the fossil record]
The fossil record is \cle{incomplete}, and you must say so in an essay. Reasons:
\begin{itemize}
\item Fossilisation is rare: most organisms decay or are eaten before burial.
\item The record is biased towards \cle{hard-bodied} organisms (bones, shells) and towards aquatic habitats where sediment accumulates. Soft-bodied worms and jellyfish fossilise very poorly.
\item Many fossils remain buried and undiscovered; erosion destroys others.
\end{itemize}
Gaps in the record therefore do \emph{not} disprove evolution; they reflect the rarity of fossilisation. Never claim the fossil record is ``complete proof'' — the examiner expects this nuance.
\end{attention}

\courssub{Comparative Anatomy: Homologous Structures}

\textbf{Observation first.} Look at your own arm, then imagine the wing of a bat, the flipper of a whale and the wing of a bird. Used for grasping, flying and swimming — utterly different functions. Yet if you X-ray them, the same bones appear in the same order: one upper-arm bone (humerus), two forearm bones (radius and ulna), a group of wrist bones (carpals), and five digits. Why would a designer reuse the same plan for such different jobs? Evolution answers: because all these animals inherited the plan from a common ancestor.

\begin{definition}[Homologous structure]
\cle{Homologous structures} are structures in different species that have the same basic plan and embryonic origin because they were inherited from a common ancestor, even though they may now perform different functions. The modification of a common ancestral plan for different functions is called \cle{divergent evolution} (adaptive radiation at a larger scale).
\end{definition}

The classic example is the \cle{pentadactyl limb} (five-digit limb) of vertebrates.

\begin{popfigure}
\begin{center}
\begin{tikzpicture}[scale=0.9, every node/.style={font=\small}]
% human arm
\begin{scope}[shift={(0,0)}]
\draw[thick, popblue, fill=popblueL] (0,0) -- (1.6,0.25) -- (1.6,-0.25) -- cycle;
\draw[thick, poporange, fill=popgold!40] (1.7,0.15) -- (3.1,0.35) -- (3.1,0.1) -- cycle;
\draw[thick, poporange, fill=popgold!40] (1.7,-0.15) -- (3.1,-0.1) -- (3.1,-0.35) -- cycle;
\draw[thick, popgreen, fill=popgreenL] (3.2,0) rectangle (3.7,0.3);
\draw[thick, popgreen, fill=popgreenL] (3.2,-0.3) rectangle (3.7,0);
\foreach \y in {-0.25,-0.12,0.01,0.14,0.27}{\draw[thick, poppurple] (3.8,\y) -- (4.7,\y+0.06);}
\node[below] at (2.3,-0.6) {\textbf{Human} (grasping)};
\end{scope}
% bat wing
\begin{scope}[shift={(6.5,0)}]
\draw[thick, popblue, fill=popblueL] (0,0) -- (1.2,0.2) -- (1.2,-0.2) -- cycle;
\draw[thick, poporange, fill=popgold!40] (1.3,0.1) -- (2.3,0.25) -- (2.3,0) -- cycle;
\draw[thick, poporange, fill=popgold!40] (1.3,-0.1) -- (2.3,-0.05) -- (2.3,-0.3) -- cycle;
\draw[thick, popgreen, fill=popgreenL] (2.4,-0.1) rectangle (2.7,0.2);
\draw[thick, poppurple] (2.75,0.15) -- (4.3,0.9);
\draw[thick, poppurple] (2.75,0.05) -- (4.3,0.35);
\draw[thick, poppurple] (2.75,-0.02) -- (4.2,-0.2);
\draw[thick, poppurple] (2.75,-0.08) -- (4.0,-0.7);
\draw[thick, poppurple] (2.7,-0.05) -- (3.2,-0.9);
\node[below] at (2.1,-1.1) {\textbf{Bat} (flying)};
\end{scope}
% whale flipper
\begin{scope}[shift={(0,-3.2)}]
\draw[thick, popblue, fill=popblueL] (0,0) -- (1.4,0.3) -- (1.4,-0.3) -- cycle;
\draw[thick, poporange, fill=popgold!40] (1.5,0.2) -- (2.6,0.4) -- (2.6,0.1) -- cycle;
\draw[thick, poporange, fill=popgold!40] (1.5,-0.2) -- (2.6,-0.1) -- (2.6,-0.4) -- cycle;
\draw[thick, popgreen, fill=popgreenL] (2.7,-0.15) rectangle (3.2,0.25);
\foreach \y in {-0.1,0.0,0.1,0.2}{\draw[thick, poppurple] (3.25,\y) -- (4.2,\y+0.1);}
\draw[thick, popink] (-0.2,0.55) .. controls (2,0.95) and (3.6,0.8) .. (4.6,0.35);
\draw[thick, popink] (-0.2,-0.55) .. controls (2,-0.95) and (3.6,-0.7) .. (4.6,-0.1);
\node[below] at (2.3,-1.1) {\textbf{Whale} (swimming)};
\end{scope}
% bird wing
\begin{scope}[shift={(6.5,-3.2)}]
\draw[thick, popblue, fill=popblueL] (0,0) -- (1.4,0.25) -- (1.4,-0.25) -- cycle;
\draw[thick, poporange, fill=popgold!40] (1.5,0.15) -- (2.7,0.3) -- (2.7,0.05) -- cycle;
\draw[thick, poporange, fill=popgold!40] (1.5,-0.15) -- (2.7,-0.05) -- (2.7,-0.3) -- cycle;
\draw[thick, popgreen, fill=popgreenL] (2.8,-0.1) rectangle (3.3,0.2);
\draw[thick, poppurple] (3.35,0.12) -- (4.4,0.35);
\draw[thick, poppurple] (3.35,0.0) -- (4.3,0.0);
\draw[thick, poppurple] (3.35,-0.06) -- (4.1,-0.35);
\node[below] at (2.2,-1.1) {\textbf{Bird} (flying)};
\end{scope}
% legend
\node[anchor=west] at (0,-5.2) {\textcolor{popblue}{\rule{0.4cm}{0.18cm}} humerus \quad \textcolor{poporange}{\rule{0.4cm}{0.18cm}} radius \& ulna \quad \textcolor{popgreen}{\rule{0.4cm}{0.18cm}} carpals \quad \textcolor{poppurple}{\rule{0.4cm}{0.18cm}} digits};
\end{tikzpicture}
\end{center}
\legende{The pentadactyl limb in four vertebrates. The same bones (same colours) occur in the same positions despite very different functions — evidence of inheritance from a common ancestor, modified by divergent evolution.}
\end{popfigure}

\begin{attention}[Homologous vs analogous — the classic trap]
Do \textbf{not} cite the wing of a bird and the wing of an insect as evidence of common ancestry. These are \cle{analogous structures}: they perform the same function (flight) but have completely different basic plans and different embryonic origins. They result from \cle{convergent evolution} — similar selection pressures producing similar solutions in unrelated groups. Only \emph{homologous} structures (same plan, different function) demonstrate common ancestry. Writing ``bird wings and insect wings show common descent'' will cost you marks.
\end{attention}

\textbf{Vestigial organs.} A related argument comes from structures that have lost their function.

\begin{definition}[Vestigial organ]
A \cle{vestigial organ} is a reduced, functionless (or much-reduced) structure in an organism, which was fully developed and functional in its ancestors. Its existence is evidence of evolutionary history.
\end{definition}

Examples: the human \cle{appendix} (a remnant of the larger caecum used by herbivorous ancestors to digest cellulose), the \cle{pelvic bones of whales and pythons} (remnants of hind limbs from walking ancestors), the muscles that move the human ear, and the wings of flightless birds such as the ostrich. Why would a whale carry hip bones it does not use, unless its ancestors walked on land?

\courssub{Comparative Embryology}

\textbf{Observation first.} If you were shown early embryos of a fish, a salamander, a chicken and a human, you could not tell them apart. All pass through stages with \cle{pharyngeal (gill) slits} and a post-anal tail — even humans, in whom the slits become parts of the ear and throat and the tail is reduced to the coccyx.

The evolutionary explanation is direct: early vertebrate embryos resemble one another because all vertebrates share a common ancestor, and the early developmental programme has been conserved. Development is modified mainly at later stages, so the shared ancestry is most visible early on.

\begin{attention}
Avoid the outdated slogan ``ontogeny recapitulates phylogeny'' (the claim that an embryo replays its whole evolutionary history). Modern biology says only that early embryos of related groups are strikingly \emph{similar}, reflecting shared developmental genes. State it that way.
\end{attention}

\courssub{Molecular Evidence}

\textbf{Observation first.} Every living cell uses the same molecule — DNA — with the same genetic code, the same 20 amino acids, and near-identical machinery for protein synthesis. This universality is itself powerful evidence that all life descends from a common ancestor. But molecular biology goes further: it allows us to \emph{measure} relatedness.

The logic is simple. Mutations accumulate in DNA (and hence in proteins) at an approximately steady rate over long periods. Two species that diverged recently have had little time to accumulate differences; two species that diverged long ago have many. The number of sequence differences is therefore a \cle{molecular clock}.

\begin{propriete}[The molecular clock principle]
For a given gene or protein, the number of differences in DNA or amino acid sequence between two species is roughly proportional to the time since they diverged from their common ancestor. The more similar the sequences, the more recent the common ancestor. (The reasoning is examinable; no calculation of dates is required at this level.)
\end{propriete}

\begin{exemplebox}[Cytochrome c and haemoglobin]
\emph{Cytochrome c} is a respiratory protein found in all aerobic organisms. Humans and chimpanzees have \emph{identical} cytochrome c; rhesus monkeys differ from humans at about 1 amino acid; dogs at about 11; chickens at about 13; yeast at more than 40. The same pattern appears for the \emph{haemoglobin} chains: human and chimpanzee beta-chains are identical, while fish haemoglobin differs at dozens of positions. The molecular family tree matches the tree built from fossils and anatomy — independent evidence converging on the same conclusion.
\end{exemplebox}

\begin{popfigure}
\begin{center}
\begin{tikzpicture}
\begin{axis}[
    ybar, bar width=0.7cm,
    width=11cm, height=6cm,
    ymin=0, ymax=80,
    ylabel={Amino acid differences from human haemoglobin},
    symbolic x coords={Chimpanzee, Chicken, Fish},
    xtick=data,
    nodes near coords,
    ymajorgrids=true,
    grid style={popline, dashed},
    every node near coord/.append style={font=\small\bfseries, color=popdark},
]
\addplot[fill=popblue, draw=popdark] coordinates {(Chimpanzee,0) (Chicken,35) (Fish,68)};
\end{axis}
\end{tikzpicture}
\end{center}
\legende{Approximate number of amino acid differences between human beta-haemoglobin and that of other vertebrates (values illustrative of the real pattern). Differences increase with time since the common ancestor — the molecular clock.}
\end{popfigure}

\begin{savaistu}[Exam tip — the strongest evidence]
Molecular evidence is now regarded as the \cle{strongest} evidence for evolution, because it is quantitative, objective and independent of anatomy. In an essay, quote a concrete example: identical \emph{cytochrome c} in humans and chimpanzees, DNA hybridisation showing the degree of DNA pairing between species, or the universality of the genetic code. A named, explained molecular example earns more credit than vague statements about ``DNA''.
\end{savaistu}

\courssub{Biogeography}

\textbf{Observation first.} Why does Australia have kangaroos, koalas and wombats — marsupials found almost nowhere else — while placental mammals dominate Africa and Eurasia? Why are the finches of the Gal\'apagos and the giant tortoises of island chains similar to mainland species, yet distinct?

\cle{Biogeography}, the study of the geographical distribution of species, makes sense only in the light of evolution combined with isolation (Section 3):

\begin{itemize}
\item \textbf{Island endemics.} Islands are colonised by a few founders from the nearest mainland; isolated and subject to different selection pressures, they evolve into new species found nowhere else (\cle{endemic} species). Darwin's finches are the classic case; the unique flora of Mount Cameroon and the endemic chameleons of the Cameroon highlands illustrate the same principle.
\item \textbf{Australian marsupials.} Australia separated from the other continents before placental mammals became dominant. Its marsupials then underwent their own adaptive radiation, producing marsupial equivalents of wolves, mice and moles — a spectacular case of divergent evolution in isolation, with convergence on placental ``counterparts'' elsewhere.
\item \textbf{Continental drift.} The distribution of fossils and living groups matches the break-up of the ancient supercontinent Gondwana, explaining why related fossils occur on now-distant continents such as Africa and South America.
\end{itemize}

\courssub{Direct Observation of Evolution}

Evolution is not only history written in rocks and molecules; it can be \emph{watched} happening today.

\begin{itemize}
\item \textbf{Antibiotic resistance.} When bacteria (e.g.\ the tuberculosis bacillus, a real concern in Cameroonian clinics) are exposed to an antibiotic, most die, but rare resistant mutants survive and multiply. Within a few bacterial generations the population is resistant — natural selection observed in real time. This is why patients must complete their full course of antibiotics.
\item \textbf{Industrial melanism.} The peppered moth \emph{Biston betularia} was mostly light-coloured before industrialisation, camouflaged on lichen-covered trees. As soot darkened tree trunks in industrial England, the dark (melanic) mutant form was better camouflaged, survived predation by birds, and became common; after clean-air laws, the light form recovered. Selection in action, documented within a human lifetime.
\item \textbf{Laboratory selection experiments.} Selection on \emph{Drosophila} (fruit flies) for bristle number, or on maize for oil content, produces large, rapid changes over tens of generations, demonstrating that selection can shift a population far beyond its original range of variation.
\end{itemize}

\begin{exoresolu}[Worked essay extract: evidence for evolution]
\textbf{Question (GCE style):} Describe four lines of evidence that support the theory of organic evolution. \hfill (8 marks)
\tcblower
\textbf{Model answer (one line of evidence = one explained example):}
\begin{enumerate}
\item \textbf{Fossil evidence:} fossils arranged in strata show a succession from simple to complex forms; \emph{Archaeopteryx} combines reptilian teeth and tail with bird feathers, showing a transition from reptiles to birds.
\item \textbf{Comparative anatomy:} the pentadactyl limb of human, bat, whale and bird has the same bones in the same arrangement despite different functions — homologous structures inherited from a common ancestor and modified by divergent evolution.
\item \textbf{Molecular evidence:} the cytochrome c of humans and chimpanzees is identical, while more distantly related species show more amino acid differences, as expected if differences accumulate with time since divergence (molecular clock).
\item \textbf{Direct observation:} bacteria exposed to antibiotics evolve resistance within a few generations, and the peppered moth changed from light to dark forms during industrial pollution — natural selection observed directly.
\end{enumerate}
Each point names the line of evidence, gives a \emph{specific example} and \emph{explains why} it supports evolution — the three ingredients of full marks.
\end{exoresolu}

\begin{methode}[Answering ``describe the evidence for evolution'' essays]
\begin{enumerate}
\item Choose three to five \emph{distinct} lines of evidence (fossils, anatomy, embryology, molecular, biogeography, direct observation).
\item For each line, follow the formula: \textbf{name it $\rightarrow$ give one precise example $\rightarrow$ explain the evolutionary logic}.
\item Include at least one molecular example (cytochrome c, haemoglobin, DNA hybridisation) — it is considered the strongest evidence.
\item Avoid analogous structures and avoid claiming the fossil record is complete; mention its bias towards hard-bodied organisms if you discuss fossils critically.
\item Conclude by noting that the \emph{convergence} of independent lines of evidence makes the case overwhelming.
\end{enumerate}
\end{methode}

\begin{aretenir}[Key points of Section 4]
\begin{itemize}
\item Fossils form by rapid burial and mineralisation; the fossil record shows an ordered succession of forms and includes transitional forms such as \emph{Archaeopteryx} and the horse series — but it is incomplete and biased towards hard-bodied organisms.
\item Homologous structures (pentadactyl limb) show common ancestry through divergent evolution; analogous structures (bird and insect wings) show convergence and are \emph{not} evidence of common descent.
\item Vestigial organs (human appendix, whale pelvis) are remnants of functional ancestral structures.
\item Early vertebrate embryos are strikingly similar, reflecting shared developmental programmes.
\item Molecular comparisons (cytochrome c, haemoglobin, DNA) quantify relatedness via the molecular clock and are now the strongest evidence.
\item Biogeography (island endemics, Australian marsupials) and direct observations (antibiotic resistance, industrial melanism, selection experiments) confirm that evolution is happening still.
\end{itemize}
\end{aretenir}

\begin{exercice}[Check your understanding]
\begin{enumerate}
\item Define the term \emph{fossil} and describe the main steps by which a fossil of a hard-bodied animal forms.
\item Explain why \emph{Archaeopteryx} is described as a transitional form.
\item Distinguish between homologous and analogous structures, giving one example of each, and state which provides evidence of common ancestry.
\item Human and chimpanzee cytochrome c are identical, while human and yeast cytochrome c differ at more than 40 positions. What does this indicate about their evolutionary relationships, and what principle does it illustrate?
\item Give two reasons why the fossil record is incomplete, and explain why this does not disprove evolution.
\end{enumerate}
\end{exercice}

\begin{corrige}[Exercise 1]
\begin{enumerate}
\item A fossil is any preserved remains, impression or trace of an organism from the geological past. Steps: death and sinking into sediment; decay of soft parts; burial of hard parts by sediment layers; compaction into rock and petrification by mineral replacement; exposure by erosion or excavation.
\item It combines features of two major groups: feathers and wings (birds) with teeth, clawed wing fingers and a long bony tail (reptiles), showing an intermediate stage in the evolution of birds from reptiles.
\item Homologous: same basic plan and origin, different functions — e.g.\ the pentadactyl limb of human and bat; evidence of common ancestry. Analogous: same function, different plan and origin — e.g.\ wings of birds and insects; evidence only of convergent evolution, not common descent.
\item Humans and chimpanzees share a very recent common ancestor (no time for differences to accumulate); humans and yeast diverged far earlier. This illustrates the molecular clock: sequence differences accumulate roughly in proportion to time since divergence.
\item Fossilisation is a rare event requiring rapid burial; the record is biased towards hard-bodied organisms and aquatic sediments; many fossils are destroyed by erosion or remain undiscovered. Incompleteness reflects the rarity of preservation, not the absence of evolution — the many transitional forms that \emph{are} known still confirm the pattern.
\end{enumerate}
\end{corrige}

With the evidence assembled, one question remains: \emph{how} did scientists explain the mechanism of change? That is the story of Lamarck, Darwin and the modern synthesis — the subject of Section 5, where we will also consolidate the whole chapter for the GCE examination.

\coursec{Section 5: Theories of Evolution and GCE Revision}

In Section 4 we examined the \cle{evidence} for organic evolution: fossils, comparative anatomy, embryology, biogeography and molecular biology all point to the same conclusion — species change over time and share common ancestors. But evidence alone does not explain \emph{how} evolution happens. For more than two centuries, biologists have proposed \cle{theories of evolution} to account for the mechanism. In this final section we compare the great historical theories — Lamarckism, Darwin--Wallace natural selection, and the modern synthetic theory — we examine the tempo of evolution, we outline human evolution, and we close with a complete revision strategy for the GCE Advanced Level examination.

\courssub{Lamarck's Theory: Use, Disuse and Acquired Characteristics}

Jean-Baptiste de Lamarck (1744--1829) was the first scientist to propose a coherent, testable mechanism for evolution. His theory, published in 1809 in his \emph{Philosophie zoologique}, rests on two laws.

\begin{definition}[Lamarck's theory of evolution]
Lamarck proposed that:
\begin{enumerate}
\item \textbf{Law of use and disuse:} organs that are used intensively become larger and more developed, while organs that are not used degenerate and may disappear.
\item \textbf{Inheritance of acquired characteristics:} modifications acquired by an organism during its lifetime are transmitted to its offspring, so that changes accumulate over generations and transform the species.
\end{enumerate}
\end{definition}

\begin{exemplebox}[The giraffe's neck according to Lamarck]
Lamarck explained the giraffe's long neck as follows: ancestral giraffes had short necks; as they stretched continually to reach leaves high in the trees, their necks lengthened slightly during each individual's lifetime; this acquired lengthening was passed to offspring; and after many generations of stretching, the modern long-necked giraffe appeared.
\end{exemplebox}

\textbf{Why Lamarck's theory is rejected.} The theory fails on both experimental and theoretical grounds:

\begin{itemize}
\item \textbf{Acquired characteristics are not inherited.} A bodybuilder's children are not born muscular; the circumcision of boys, practised for thousands of generations in many communities, has never produced boys born without a foreskin. Changes to body (somatic) cells do not alter the DNA of gametes.
\item \textbf{Weismann's experiment.} August Weismann cut off the tails of mice for many successive generations; every new generation was still born with full-length tails. Mutilations acquired during life are not transmitted.
\item \textbf{Modern genetics explains why.} Only changes in the \cle{genotype} of germ-line cells (mutations in gametes or their precursors) can be inherited. The environment can modify the \emph{phenotype} but not the hereditary material of the gametes.
\end{itemize}

\begin{attention}[Common mistake: presenting Lamarck as accepted science]
A very frequent GCE error is to write, when explaining natural selection, that giraffes ``stretched their necks and passed longer necks to their offspring''. This is \textbf{Lamarckism}, not Darwinism, and it will cost you marks. In Darwin's correct explanation, long-necked variants \emph{already existed} in the population by random variation; they survived and reproduced \emph{better}; they did not acquire long necks by effort. Never describe evolution as the result of an organism's ``need'', ``will'' or ``effort''.
\end{attention}

\courssub{The Darwin--Wallace Theory of Natural Selection}

Charles Darwin and Alfred Russel Wallace independently arrived at the same mechanism. Wallace, working in the Malay Archipelago, sent Darwin an essay outlining the idea, and their work was presented jointly to the Linnean Society of London in 1858; Darwin then developed the theory fully in \emph{On the Origin of Species} (1859). The theory follows logically from three observations and two deductions:

\begin{enumerate}
\item \textbf{Observation 1:} organisms produce far more offspring than can possibly survive (overproduction).
\item \textbf{Observation 2:} despite this, population sizes remain roughly stable, so there is a \cle{struggle for existence} (competition for food, mates, space; predation; disease).
\item \textbf{Observation 3:} individuals within a species show heritable \cle{variation} — no two individuals are identical.
\item \textbf{Deduction 1:} individuals whose variations make them better adapted survive and reproduce more successfully: this is \cle{natural selection}, often summarised as ``survival of the fittest''.
\item \textbf{Deduction 2:} favourable variations accumulate over many generations, so the population gradually changes and may eventually form a new species.
\end{enumerate}

\begin{definition}[Fitness]
In evolutionary biology, \cle{fitness} means \emph{reproductive success}: the contribution an individual makes to the gene pool of the next generation, relative to other individuals in the population. It does not mean physical strength. A small, weak organism that leaves many surviving offspring is fitter than a strong one that leaves none.
\end{definition}

\textbf{Strengths of the theory.} It is mechanistic (no ``will'' or purpose needed), it explains adaptation, it explains the unity and diversity of life from common ancestors, and it makes testable predictions confirmed by countless observations — industrial melanism in the peppered moth, antibiotic resistance in bacteria such as \emph{Mycobacterium tuberculosis}, and insecticide resistance in \emph{Anopheles} mosquitoes here in Cameroon.

\textbf{Original weaknesses.} Darwin could not explain:
\begin{itemize}
\item the \emph{source} of variation (he knew nothing of mutation);
\item the \emph{mechanism of inheritance} (he wrongly imagined blending inheritance, which would dilute favourable variations);
\item why variation does not simply disappear by blending over generations.
\end{itemize}
Mendel's work (1866), unknown to Darwin, would later supply the missing answers.

\courssub{The Modern Synthetic Theory (Neo-Darwinism)}

\begin{definition}[Modern synthetic theory of evolution]
The \cle{modern synthetic theory} (or neo-Darwinism) is the fusion of Darwin--Wallace natural selection with Mendelian genetics and mutation theory. It states that evolution is the change in \cle{allele frequencies} within a population's \cle{gene pool} over generations, produced by mutation (the ultimate source of new alleles), genetic recombination, natural selection, genetic drift, and isolation. Natural selection acts on phenotypes, but what is transmitted and what evolves is the genetic composition of the population.
\end{definition}

The synthesis resolves Darwin's weaknesses elegantly:
\begin{itemize}
\item \textbf{Source of variation:} random \cle{mutations} create new alleles; meiosis (crossing over and independent assortment) and random fertilisation reshuffle them into new combinations.
\item \textbf{Inheritance:} genes are particulate, not blending, so favourable alleles are preserved intact and can increase in frequency.
\item \textbf{Selection:} the environment ``screens'' phenotypes; individuals with advantageous genotypes contribute more alleles to the next generation's gene pool.
\end{itemize}

\begin{popfigure}
\begin{center}
\begin{tikzpicture}[scale=0.92, every node/.style={font=\small}]
% Lamarck column
\draw[fill=popblueL, draw=popblue] (-0.3,6.2) rectangle (3.9,7.2);
\node at (1.8,6.7) {\textbf{Lamarck}};
\draw[fill=popgreenL, draw=popgreen] (4.3,6.2) rectangle (8.5,7.2);
\node at (6.4,6.7) {\textbf{Darwin--Wallace}};
\draw[fill=popblueL, draw=poppurple] (8.9,6.2) rectangle (13.1,7.2);
\node at (11,6.7) {\textbf{Neo-Darwinism}};
% row 1: mechanism
\node[align=center, text width=3.6cm] at (1.8,5.2) {Giraffe stretches to reach leaves; neck lengthens by effort};
\node[align=center, text width=3.6cm] at (6.4,5.2) {Neck length varies naturally; long-necked giraffes survive better};
\node[align=center, text width=3.6cm] at (11,5.2) {Alleles for longer necks arise by mutation and increase in frequency};
% row 2: inheritance
\node[align=center, text width=3.6cm] at (1.8,3.4) {Acquired long neck is inherited by offspring};
\node[align=center, text width=3.6cm] at (6.4,3.4) {Favourable variants are inherited (mechanism unknown to Darwin)};
\node[align=center, text width=3.6cm] at (11,3.4) {Selected alleles are transmitted via gametes (Mendelian genetics)};
% row 3: verdict
\node[align=center, text width=3.6cm, text=popdark] at (1.8,1.6) {\textbf{Rejected:} acquired traits are not inherited};
\node[align=center, text width=3.6cm, text=popdark] at (6.4,1.6) {\textbf{Correct mechanism,} but no genetics};
\node[align=center, text width=3.6cm, text=popdark] at (11,1.6) {\textbf{Accepted:} selection + genetics + mutation};
% row labels
\node[align=right, font=\small\bfseries, text width=2.2cm] at (-1.6,5.2) {Mechanism};
\node[align=right, font=\small\bfseries, text width=2.2cm] at (-1.6,3.4) {Inheritance};
\node[align=right, font=\small\bfseries, text width=2.2cm] at (-1.6,1.6) {Verdict};
\draw[popline] (-0.3,0.8) -- (13.1,0.8);
\end{tikzpicture}
\end{center}
\legende{Comparison of the three theories of evolution using the giraffe's neck.}
\end{popfigure}

\courssub{Gradualism versus Punctuated Equilibrium}

Within the modern synthesis, two models describe the \emph{tempo} of evolution. \cle{Phyletic gradualism} holds that species change slowly and continuously by the steady accumulation of small variations. \cle{Punctuated equilibrium} (Eldredge and Gould, 1972) holds that species remain stable (in stasis) for long periods, interrupted by relatively rapid bursts of change, often associated with speciation in small isolated populations. The fossil record, with its long periods of stability and sudden appearances of new forms, is often better fitted by punctuated equilibrium. The two models are not rivals to natural selection — both are driven by it; they differ only in rate.

\begin{popfigure}
\begin{center}
\begin{tikzpicture}
\begin{axis}[
  width=12cm, height=7cm,
  xlabel={Time (millions of years)}, ylabel={Morphological change},
  xmin=0, xmax=10, ymin=0, ymax=10,
  xtick=\empty, ytick=\empty,
  legend style={at={(0.02,0.98)}, anchor=north west, font=\small},
  axis lines=left]
% gradualism: smooth diagonal
\addplot[popblue, very thick, smooth] coordinates {(0,0.5) (2,2.2) (4,4.0) (6,5.8) (8,7.6) (10,9.4)};
\addlegendentry{Gradualism}
% punctuated equilibrium: steps
\addplot[poppink, very thick] coordinates {(0,0.8) (3.5,0.8) (3.5,3.2) (6.5,3.2) (6.5,6.2) (8.8,6.2) (8.8,9.2) (10,9.2)};
\addlegendentry{Punctuated equilibrium}
\end{axis}
\end{tikzpicture}
\end{center}
\legende{Gradualism (smooth, continuous change) versus punctuated equilibrium (long stasis interrupted by rapid change).}
\end{popfigure}

\begin{savaistu}[Exam-useful extra]
Questions on the \emph{tempo} of evolution occasionally appear as part (c) of an essay. A single well-labelled sketch like the one above, plus one sentence on each model, usually earns full marks for that part.
\end{savaistu}

\courssub{An Outline of Human Evolution}

Humans belong to the family Hominidae. The fossil record, overwhelmingly African, documents a transition from ape-like ancestors to modern humans over roughly 6--7 million years. Note that evolution here is a \emph{bush}, not a ladder: several hominin species coexisted, and \emph{Homo sapiens} did not descend in a straight line from the forms listed below — they are best regarded as close relatives illustrating the major trends.

\begin{itemize}
\item \textbf{\emph{Australopithecus}} (about 4--2 million years ago, e.g.\ \emph{A. afarensis}, the famous ``Lucy'' from Ethiopia): small brain (about $400\ \mathrm{cm^3}$), bipedal locomotion but still adapted to climbing, ape-like face with projecting jaws.
\item \textbf{\emph{Homo habilis}} (about 2.4--1.6 million years ago, East Africa): larger brain (about $650\ \mathrm{cm^3}$), associated with the first simple stone tools (``handy man'').
\item \textbf{\emph{Homo erectus}} (about 1.9 million--100\,000 years ago): brain about $900\ \mathrm{cm^3}$, fully bipedal, used fire, made more advanced tools, first hominin to leave Africa.
\item \textbf{\emph{Homo sapiens}} (from about 300\,000 years ago, Africa): brain about $1350\ \mathrm{cm^3}$, high forehead, reduced brow ridges, prominent chin, complex language, art and culture; spread across the whole world.
\end{itemize}

Key trends in hominin evolution: increasing \cle{brain capacity}, increasingly efficient \cle{bipedalism} (foramen magnum moving forward, S-shaped spine, bowl-shaped pelvis), flattening of the face, reduction of teeth and jaws, and increasing reliance on tools and culture.

\begin{popfigure}
\begin{center}
\begin{tikzpicture}[scale=0.95, every node/.style={font=\small}]
% timeline axis
\draw[->, thick, popline] (0,0) -- (0,7.4) node[above] {Millions of years ago};
\foreach \y/\t in {0.5/0, 2/2, 3.5/4, 5/6} {
  \draw[popline] (-0.15,\y) -- (0.15,\y);
  \node[left] at (-0.2,\y) {\t};
}
% branches
\draw[very thick, popgreen] (1.2,0.4) -- (1.2,2.6) node[right, text=popdark] {\emph{Australopithecus}};
\draw[very thick, popblue] (2.6,0.4) -- (2.6,2.1) node[right, text=popdark] {\emph{Homo habilis}};
\draw[very thick, poporange] (4.0,0.4) -- (4.0,1.8) node[right, text=popdark] {\emph{Homo erectus}};
\draw[very thick, poppurple] (5.4,0.1) -- (5.4,0.9) node[right, text=popdark] {\emph{Homo sapiens}};
% connecting trunk
\draw[thick, popmuted, dashed] (1.2,2.6) -- (2.6,2.1);
\draw[thick, popmuted, dashed] (2.6,2.1) -- (4.0,1.8);
\draw[thick, popmuted, dashed] (4.0,1.8) -- (5.4,0.9);
% brain sizes
\node[align=left, text width=4.6cm, font=\small] at (9.6,2.3) {$\approx 400\ \mathrm{cm^3}$ brain; bipedal, still climbs};
\node[align=left, text width=4.6cm, font=\small] at (9.6,1.5) {$\approx 650\ \mathrm{cm^3}$; first stone tools};
\node[align=left, text width=4.6cm, font=\small] at (9.6,0.85) {$\approx 900\ \mathrm{cm^3}$; fire; leaves Africa};
\node[align=left, text width=4.6cm, font=\small] at (9.6,0.2) {$\approx 1350\ \mathrm{cm^3}$; language, art, culture};
\draw[popline, thin] (5.6,2.2) -- (7.3,2.3);
\draw[popline, thin] (5.6,1.5) -- (7.3,1.5);
\draw[popline, thin] (5.6,0.9) -- (7.3,0.85);
\draw[popline, thin] (5.6,0.45) -- (7.3,0.2);
\end{tikzpicture}
\end{center}
\legende{Simplified outline of human evolution. Note the increasing brain size and the African origin of all four groups; the dashed links indicate broad relationships, not a strict linear descent.}
\end{popfigure}

\begin{savaistu}[Africa, cradle of humanity]
The oldest known fossils of \emph{Homo sapiens} and of all earlier hominins come from Africa — Ethiopia, Kenya, Tanzania, South Africa and Chad. Genetic evidence (mitochondrial DNA diversity is greatest in African populations) confirms that every human alive today descends from African ancestors. Cameroon itself lies within the range of our closest living relatives, the chimpanzee and the gorilla.
\end{savaistu}

\courssub{Worked Example: Comparing Theories in an Examination Context}

\begin{exoresolu}[Explaining antibiotic resistance using two theories]
Tuberculosis bacteria (\emph{Mycobacterium tuberculosis}) sometimes become resistant to the antibiotic rifampicin. Explain the origin of this resistance (a) according to Lamarck's theory, (b) according to the modern synthetic theory, and state which explanation is correct and why.
\tcblower
\textbf{(a) Lamarckian explanation (incorrect):} the bacteria, ``needing'' to survive the antibiotic, would develop resistance during their lifetime and transmit this acquired resistance to their descendants.

\textbf{(b) Modern synthetic explanation (correct):} within any large bacterial population, random \cle{mutations} occur continuously; by chance, one bacterium may carry a mutation (e.g.\ in the \emph{rpoB} gene) that confers resistance \emph{before} any antibiotic is present. When rifampicin is applied, sensitive bacteria are killed (selection pressure), while the resistant mutant survives and reproduces. Its resistant allele is passed to all its offspring by normal inheritance of DNA, so the frequency of the resistance allele in the population rises rapidly. The population has \emph{evolved} resistance.

\textbf{Why (b) is correct:} experiments (e.g.\ replica plating) show that resistant mutants exist in populations that have never encountered the antibiotic; the antibiotic does not create the mutation, it only \emph{selects} it. Acquired characteristics are not inherited, so Lamarck's explanation fails. This is why patients in Cameroon and worldwide must complete their full course of anti-TB drugs: incomplete treatment selects resistant bacteria.
\end{exoresolu}

\courssub{GCE Revision Strategy for Evolution}

\textbf{Past-paper analysis.} Evolution is one of the most reliably examined topics in the GCE Advanced Level Biology papers:

\begin{itemize}
\item \textbf{Paper 1 (multiple choice):} expect short items on definitions (natural selection, speciation, gene pool), on identifying Lamarckian versus Darwinian statements, on types of selection, and on evidence of evolution. Speed and precision with vocabulary are essential.
\item \textbf{Paper 2 (structured questions):} data-based questions showing graphs of allele frequencies, tables of fossil measurements, or the peppered moth / sickle cell examples, asking you to describe trends, explain them by natural selection, and draw conclusions.
\item \textbf{Paper 2 (essay):} classic 20-mark essays include ``Describe the theories of evolution put forward by Lamarck and Darwin, and explain how the modern synthetic theory improves on both'', and ``Discuss the evidence for organic evolution''.
\end{itemize}

\begin{methode}[Structuring a 20-mark GCE essay on theories of evolution]
\begin{enumerate}
\item \textbf{Define} evolution briefly in the introduction (change in allele frequencies / descent with modification). (2 marks)
\item \textbf{Lamarck:} state his two laws, give the giraffe example, then give two reasons for rejection (Weismann's mice; genetics of soma versus germ line). (5 marks)
\item \textbf{Darwin--Wallace:} present the three observations and two deductions logically; give one supporting example (peppered moth or antibiotic resistance); state the two weaknesses (no source of variation, no mechanism of inheritance). (6 marks)
\item \textbf{Modern synthesis:} define it as the fusion of selection with Mendelian genetics and mutation; show explicitly how it cures each of Darwin's weaknesses. (5 marks)
\item \textbf{Conclusion:} one or two sentences summarising why the synthetic theory is accepted today. (2 marks)
\end{enumerate}
Write in clear paragraphs, use correct terminology, and include one example per theory. Examiners award marks for \emph{logical structure}, not for length alone.
\end{methode}

\courssub{Preparing the Practical Examination}

Practical and data-handling questions on evolution typically present a table or graph — for example, allele frequencies of the sickle-cell allele in regions of differing malaria prevalence, or moth counts before and after industrial pollution — and ask you to interpret them.

\begin{attention}[Exam tip: always conclude from the data]
In practical papers, the single most common way to lose marks is to give a vague conclusion (``evolution occurred''). Always state a \textbf{clear conclusion that refers back to the actual data}: quote figures, describe the trend, and link it to a mechanism. For example: ``The frequency of the Hb$^S$ allele is highest (about 0.15) in the high-malaria region and lowest (about 0.02) where malaria is absent, supporting the conclusion that malaria acts as a selective agent maintaining the heterozygote advantage.''
\end{attention}

\begin{methode}[Interpreting an allele-frequency graph or table]
\begin{enumerate}
\item Read the axes or column headings carefully; note units and scales.
\item Describe the trend factually first (``allele A rises from 0.1 to 0.8 over 50 generations'').
\item Identify the selective agent suggested by the context (predation, disease, antibiotic, climate).
\item Explain the trend by natural selection: variation $\rightarrow$ differential survival $\rightarrow$ change in allele frequency.
\item State a conclusion that quotes the data and names the mechanism.
\end{enumerate}
\end{methode}

\begin{exoresolu}[Interpreting data on the peppered moth]
The table shows the percentage of dark (melanic) \emph{Biston betularia} moths in a polluted industrial area over time:
\begin{center}
\begin{tabular}{|l|ccccc|}
\hline
\rowcolor{popblueL} Year & 1950 & 1960 & 1970 & 1980 & 1990 \\
\hline
Melanic moths (\%) & 10 & 45 & 80 & 90 & 94 \\
\hline
\end{tabular}
\end{center}
Describe the trend and explain it by natural selection.
\tcblower
\textbf{Description:} the percentage of melanic moths rose steadily from 10\,\% in 1950 to 94\,\% in 1990.

\textbf{Explanation:} before pollution, light moths were camouflaged on pale, lichen-covered tree bark while dark moths were easily seen and eaten by birds, so the dark allele remained rare. Industrial soot killed the lichens and blackened the bark; the dark variant was now camouflaged and the light variant conspicuous. Birds therefore ate more light moths, so dark moths survived and reproduced more, and the frequency of the dark allele (and hence of melanic moths) increased generation after generation. This is directional selection driven by differential predation — evolution observed within human timescales.
\end{exoresolu}

\courssub{Timed Mock Examination and Mark-Scheme-Based Correction}

The final stage of preparation is to sit at least one full timed paper under examination conditions (no notes, strict timing, blue or black pen). Then correct it \emph{against the mark scheme}, not against your feelings:

\begin{enumerate}
\item Tick only points that match mark-scheme creditworthy statements; note where you wrote correct ideas but missed the exact keyword (e.g.\ wrote ``genes change'' instead of ``\emph{allele frequencies} change'').
\item Classify every lost mark: (a) knowledge gap, (b) misread question, (c) poor expression, (d) time pressure. Each category has a different remedy.
\item Re-write, from memory, any answer that scored below half marks.
\item Keep an error notebook; review it in the final week.
\end{enumerate}

\begin{exercice}[Mock examination practice (45 minutes)]
\begin{enumerate}
\item Define the following terms: gene pool; natural selection; speciation; genetic drift. (4 marks)
\item State Lamarck's two laws and give one reason why his theory is rejected. (4 marks)
\item Explain, using the modern synthetic theory, how a population of \emph{Anopheles} mosquitoes becomes resistant to an insecticide. (6 marks)
\item The frequency of a recessive allele in an isolated population rose from 0.05 to 0.40 in 30 generations. Suggest two possible explanations and state what additional information would help you decide between them. (6 marks)
\end{enumerate}
\end{exercice}

\begin{corrige}[Mock examination practice]
\textbf{1.} Gene pool: the total of all alleles of all genes in a population. Natural selection: differential survival and reproduction of individuals with advantageous heritable variations, changing allele frequencies. Speciation: formation of new species from an ancestral species, usually via reproductive isolation. Genetic drift: random change in allele frequencies, significant in small populations.

\textbf{2.} Law of use and disuse; inheritance of acquired characteristics. Rejected because acquired somatic modifications do not alter gamete DNA (Weismann's tail-cutting experiment; genetics).

\textbf{3.} Random mutations produce rare resistant alleles before exposure; the insecticide kills susceptible mosquitoes (selection pressure); resistant individuals survive and reproduce; the resistance allele is inherited and its frequency rises each generation; the population evolves resistance. (Not: mosquitoes ``develop'' resistance by contact.)

\textbf{4.} (i) Natural selection favouring the allele (heterozygote advantage or recessive phenotype advantage); (ii) genetic drift, since the population is small and isolated. Helpful information: population size, whether carriers of the allele have higher survival/fertility, mutation and migration rates.
\end{corrige}

\courssub{Final Revision Checklist}

\begin{aretenir}[Revision checklist: what you must master for the GCE]
\textbf{Key terms:} evolution, variation, gene pool, allele frequency, natural selection, fitness, adaptation, directional / stabilising / disruptive selection, genetic drift, speciation, reproductive isolation, modern synthetic theory, homologous and analogous structures.

\textbf{Key examples:} peppered moth (industrial melanism); antibiotic and insecticide resistance; sickle cell and malaria in Cameroon; Darwin's finches; the giraffe's neck (to contrast theories); African hominin fossils.

\textbf{Key diagrams to reproduce from memory:} the three types of natural selection (before/after curves); the Lamarck--Darwin--neo-Darwinism comparison; the human evolution outline; gradualism versus punctuated equilibrium; a labelled isolating-mechanisms flowchart.

\textbf{Key skills:} structure a 20-mark essay logically; interpret allele-frequency data and conclude by quoting figures; never attribute evolution to an organism's ``need'' or ``effort''.
\end{aretenir}

You have now completed the chapter on natural selection and evolution — from variation in populations, through the mechanism of selection and the origin of species, to the evidence and the great theories that explain it all. Revise actively, practise past questions under time pressure, and walk into the examination hall with confidence. Evolution is the thread that ties all of biology together; master it, and every other chapter becomes easier to understand. Good luck!

\newpage

\coursec{Integration activity}

\courssub{Aims of the activity}

By the end of this investigation, you should be able to:

\begin{itemize}
\item \cle{Model} a biological population and its gene pool using simple materials, and justify the analogy between each material and the biological reality it represents.
\item Apply the \cle{Hardy--Weinberg principle} to calculate allele and genotype frequencies from real data.
\item Demonstrate experimentally that a \cle{selection pressure} (here, an insecticide) changes allele frequencies over generations, even though the total population size may recover.
\item Construct and interpret a results table and a graph of \emph{allele frequency against generation number}.
\item Write a structured scientific report in GCE style (hypothesis, method, results, conclusion, limitations).
\item Relate the model to the real public-health problem of \cle{insecticide resistance} in malaria vectors in Cameroon.
\end{itemize}

\begin{attention}[Skills being assessed]
This activity mobilises the same competences as GCE Paper 2 (structured questions) and the practical examination: manipulating data, calculating frequencies, plotting and interpreting graphs, and evaluating a model. Take the report seriously --- it is excellent examination training.
\end{attention}

\courssub{Situation}

\textbf{The problem.} In the health district of Douala, malaria remains a major cause of consultation at the Bonassama District Hospital. The main vector is the mosquito \emph{Anopheles gambiae} sensu lato. Since the early 2000s, vector control has relied heavily on pyrethroid insecticides, used in long-lasting insecticidal nets (LLINs) and in indoor residual spraying. Health workers have observed that in some neighbourhoods, mosquitoes are no longer killed by the standard dose of insecticide: the population has become \cle{resistant}.

Resistance is due, in part, to a point mutation called \emph{kdr} (``knock-down resistance'') in a gene coding for a sodium channel protein in the mosquito's nervous system. We simplify the situation to a single gene with two alleles:

\begin{itemize}
\item allele $\mathbf{R}$: confers resistance to the insecticide (dominant for survival under treatment);
\item allele $\mathbf{r}$: normal, susceptible allele.
\end{itemize}

A mosquito of genotype $RR$ or $Rr$ survives spraying; a mosquito of genotype $rr$ dies with a probability fixed by the ``dose'' of insecticide.

\begin{definition}[Gene pool and allele frequency]
The \cle{gene pool} of a population is the total of all the alleles of all the genes present in the individuals of that population at a given time. The \cle{allele frequency} of an allele is the proportion of that allele among all the alleles of the gene in the population. For a gene with two alleles $R$ and $r$, we write $p = \mathrm{freq}(R)$ and $q = \mathrm{freq}(r)$, with $p + q = 1$.
\end{definition}

\begin{propriete}[Hardy--Weinberg principle]
In a large population, with random mating, and in the absence of mutation, migration, selection and genetic drift, allele frequencies remain constant from generation to generation, and genotype frequencies are given by
\[
p^2 \;+\; 2pq \;+\; q^2 \;=\; 1
\]
where $p^2 = \mathrm{freq}(RR)$, $2pq = \mathrm{freq}(Rr)$ and $q^2 = \mathrm{freq}(rr)$. The derivation (from the random union of gametes) is examinable at GCE Advanced Level. The principle is used here as a \emph{null model}: any deviation from its predictions indicates that an evolutionary agent (such as selection) is acting.
\end{propriete}

\textbf{Your mission.} You are a team of biologists seconded to the National Malaria Control Programme. Using a bead model, you must simulate what happens to the frequency of the resistance allele in a mosquito population subjected to repeated insecticide spraying, and then advise the Programme on the consequences for vector control strategy.

\textbf{Materials per group (4 students).}

\begin{itemize}
\item 100 beads: 50 red beads (allele $R$) and 50 white beads (allele $r$) --- but see Task 1: you will choose the \emph{initial} composition;
\item one opaque bag (the ``gene pool'');
\item one six-sided die (the ``insecticide'');
\item calculator, graph paper, report sheets.
\end{itemize}

\textbf{Protocol imposed by the Programme.}

\begin{enumerate}
\item Place the chosen initial beads in the bag. The bag always represents the \emph{gametes} available to form the next generation.
\item \textbf{Forming offspring:} draw beads \emph{two by two} (without looking) to form 25 offspring genotypes ($RR$, $Rr$ or $rr$). Record the numbers of each genotype. Return all beads to the bag after counting --- the allele frequencies of the gene pool are unchanged by drawing, only \emph{expressed} as genotypes.
\item \textbf{Spraying:} for each $rr$ offspring, roll the die. If the die shows 1, 2, 3, 4 or 5, the mosquito is killed; if it shows 6, it survives (the insecticide kills $5/6 \approx 83\%$ of susceptible homozygotes). $RR$ and $Rr$ individuals always survive.
\item \textbf{Reproduction:} the survivors reproduce. Compute the new allele frequencies in the surviving population (Task 4 shows how), and reconstitute the bag with 100 beads matching these new frequencies (rounded to whole beads).
\item Repeat for \textbf{6 generations}, recording all data in a table.
\end{enumerate}

\textbf{Data sheet (to be completed).}

\begin{center}
\begin{tabularx}{\linewidth}{|c|X|X|X|X|X|X|}
\hline
\rowcolor{popblueL} Generation & $RR$ formed & $Rr$ formed & $rr$ formed & $rr$ surviving & freq($R$) after selection & freq($r$) after selection \\
\hline
0 (initial) & -- & -- & -- & -- & & \\
\hline
1 & & & & & & \\
\hline
2 & & & & & & \\
\hline
3 & & & & & & \\
\hline
4 & & & & & & \\
\hline
5 & & & & & & \\
\hline
6 & & & & & & \\
\hline
\end{tabularx}
\end{center}

\courssub{Tasks}

\begin{enumerate}
\item \textbf{Setting up the gene pool.} Resistance mutations are initially rare in an untreated population. Justify the choice of an initial composition of 10 red beads ($R$) and 90 white beads ($r$). State the initial allele frequencies $p_0 = \mathrm{freq}(R)$ and $q_0 = \mathrm{freq}(r)$.

\item \textbf{Hypothesis.} Write a clear, testable hypothesis predicting how the frequency of allele $R$ will change over the six generations of spraying, and justify it in one or two sentences using the principle of natural selection.

\item \textbf{Hardy--Weinberg check.} Assuming random mating and no selection, calculate the \emph{expected} genotype frequencies in generation 1 from $p_0$ and $q_0$ using the Hardy--Weinberg equilibrium $p^2 + 2pq + q^2 = 1$. Compare these expected numbers (out of 25 offspring) with the numbers your group actually drew. Suggest one reason for any difference.

\item \textbf{Recalculating allele frequencies after selection.} Show, using your generation 1 data, how to compute the frequency of $R$ among the \emph{survivors}:
\[
p = \frac{2N_{RR} + N_{Rr}}{2N_{\text{survivors}}}
\]
Explain the meaning of the numerator and the denominator.

\item \textbf{Graph.} Plot a graph of $\mathrm{freq}(R)$ (y-axis) against generation number (x-axis), for generations 0 to 6. Describe the shape of the curve: is the increase of $R$ fastest at the beginning, in the middle, or at the end? Explain why, considering which genotypes the insecticide can ``see''.

\item \textbf{Interpretation.} Explain, in terms of phenotypes and fitness, why allele $r$ is \emph{not} completely eliminated even after several generations of spraying. (Hint: in which genotype can $r$ ``hide'' from selection?)

\item \textbf{Link to Cameroon.} Entomologists of the Organisation de Coordination pour la lutte contre les End\'emies en Afrique Centrale (OCEAC) and of the Biotechnology Centre of the University of Yaound\'e I have reported high frequencies of the \emph{kdr} resistance mutation in \emph{Anopheles gambiae} populations in Douala and Yaound\'e. Using your model, explain to the National Malaria Control Programme:
\begin{enumerate}
\item why resistance appeared and spread once pyrethroid nets were distributed on a large scale;
\item why simply \emph{increasing the dose} of the same insecticide is not a sustainable solution;
\item one strategy, suggested by your model, that could slow down the evolution of resistance (think of what happens if the selection pressure is removed or alternated).
\end{enumerate}

\item \textbf{Limitations of the model.} Identify \emph{two} ways in which the bead model differs from a real mosquito population, and state for each one whether it would make the evolution of resistance in reality \emph{faster} or \emph{slower} than in the model. Consider at least: population size, migration of mosquitoes from untreated areas, and the cost of resistance in the absence of insecticide.

\item \textbf{Conclusion.} Write a conclusion of 5--8 lines summarising what the model demonstrates about natural selection, and one recommendation for malaria vector control in Cameroon.
\end{enumerate}

\courssub{Model answer}

\begin{methode}[Task 1 -- Initial gene pool]
New resistance alleles arise by mutation, and mutation is a rare event; before any spraying, resistant mosquitoes have no advantage, so the allele stays at low frequency. We therefore start with 10 red beads and 90 white beads:
\[
p_0 = \mathrm{freq}(R) = \frac{10}{100} = 0.10 \qquad q_0 = \mathrm{freq}(r) = \frac{90}{100} = 0.90
\]
Check: $p_0 + q_0 = 1$.
\end{methode}

\textbf{Task 2 -- Hypothesis.} ``If the insecticide kills susceptible $rr$ mosquitoes every generation, then the frequency of the resistance allele $R$ will increase from generation to generation, because carriers of $R$ survive and reproduce more; the increase will be slow at first (when $R$ is rare) and faster later.''

\textbf{Task 3 -- Hardy--Weinberg expectations for generation 1.}
\[
p^2 = (0.1)^2 = 0.01; \quad 2pq = 2(0.1)(0.9) = 0.18; \quad q^2 = (0.9)^2 = 0.81
\]
Expected offspring out of 25: $RR \approx 0.25$ (i.e. 0 or 1), $Rr \approx 4.5$ (4 or 5), $rr \approx 20.25$ (about 20). A typical draw might give $RR = 0$, $Rr = 5$, $rr = 20$. Small differences from expectation are due to \emph{sampling error}: 25 draws is a small sample, so chance fluctuations are large (this models \cle{genetic drift} in small populations).

\textbf{Task 4 -- Allele frequency after selection (worked example with typical data).}

Suppose generation 1 gave $RR = 0$, $Rr = 5$, $rr = 20$. Spraying kills $5/6$ of the $rr$: suppose the die kills 17 and 3 survive. Survivors: $N_{RR} = 0$, $N_{Rr} = 5$, $N_{rr} = 3$, total $N = 8$.

The numerator $2N_{RR} + N_{Rr}$ counts all $R$ alleles (each $RR$ carries two, each $Rr$ carries one); the denominator $2N$ is the total number of alleles in the survivors (each mosquito is diploid):
\[
p_1 = \frac{2(0) + 5}{2 \times 8} = \frac{5}{16} = 0.3125 \approx 0.31; \qquad q_1 = 1 - p_1 \approx 0.69
\]
The bag is refilled with about 31 red and 69 white beads. In one generation, $\mathrm{freq}(R)$ has roughly tripled.

\textbf{Task 5 -- Graph.} A typical set of results (your values will differ slightly because of the die):

\begin{center}
\begin{tabularx}{\linewidth}{|l|X|X|X|X|X|X|X|}
\hline
\rowcolor{popblueL} Generation & 0 & 1 & 2 & 3 & 4 & 5 & 6 \\
\hline freq($R$) & 0.10 & 0.31 & 0.48 & 0.62 & 0.74 & 0.83 & 0.89 \\
\hline
\end{tabularx}
\end{center}

\begin{popfigure}
\begin{tikzpicture}
\begin{axis}[
    width=11cm, height=7cm,
    xlabel={Generation}, ylabel={Frequency of allele $R$},
    xmin=0, xmax=6, ymin=0, ymax=1,
    xtick={0,1,2,3,4,5,6},
    ytick={0,0.2,0.4,0.6,0.8,1.0},
    grid=both, grid style={popline, very thin},
    axis line style={popink},
]
\addplot[color=poporange, very thick, mark=*, mark size=2pt] coordinates {
    (0,0.10) (1,0.31) (2,0.48) (3,0.62) (4,0.74) (5,0.83) (6,0.89)
};
\end{axis}
\end{tikzpicture}
\legende{Typical result: frequency of the resistance allele $R$ rises over six generations of insecticide spraying.}
\end{popfigure}

The curve rises steeply in the \emph{middle} generations and levels off near the end. At the start, $R$ is rare and mostly hidden in heterozygotes, so few extra $R$ carriers are exposed; in the middle, many $Rr$ and $RR$ individuals survive while $rr$ are eliminated, so $R$ spreads fast; at the end, the remaining $r$ alleles are almost all in $Rr$ heterozygotes, which the insecticide cannot eliminate, so the increase slows down.

\textbf{Task 6 -- Why $r$ is never eliminated.} The insecticide acts on the \emph{phenotype}: only $rr$ individuals are susceptible. In heterozygotes $Rr$, the $r$ allele is masked by the dominant $R$ and confers no disadvantage. These ``hidden'' $r$ alleles are passed on each generation. Selection can therefore push $q$ very low but, with dominance, cannot easily remove the last copies of $r$ --- this is why a recessive susceptible allele persists at low frequency, and why susceptibility can re-emerge.

\textbf{Task 7 -- Advice to the Programme.}
\begin{enumerate}
\item Distributing pyrethroid nets on a large scale applied a strong, continuous \cle{selection pressure}: susceptible mosquitoes were killed, the rare resistant ones survived and reproduced, and their offspring inherited $R$. Generation after generation (and mosquitoes have many generations per year in Douala's climate), $\mathrm{freq}(R)$ rose exactly as in our model, until resistant mosquitoes dominated the population.
\item Increasing the dose of the \emph{same} insecticide only strengthens the same selection pressure: it may kill some heterozygotes too, but any fully resistant survivors are selected even more strongly. It also increases cost and toxicity to humans and the environment without changing the evolutionary outcome.
\item The model suggests removing or varying the pressure: \emph{rotate} insecticides with different modes of action, use mosaic spraying, or combine nets with larval-source management (draining stagnant water, \emph{Bti} larviciding). Where resistance carries a fitness cost, a pause in pyrethroid use lets $r$ recover, restoring susceptibility.
\end{enumerate}

\textbf{Task 8 -- Limitations.}
\begin{itemize}
\item \textbf{Population size:} our ``population'' of 25 offspring is tiny, so chance (drift) causes large random jumps; real mosquito populations number in the millions, where selection acts more smoothly and predictably. Resistance in reality therefore evolves \emph{more steadily}, not necessarily faster.
\item \textbf{Migration:} our bag is closed; real populations receive susceptible immigrants from untreated areas (no gene flow in the model). Immigration of $r$ alleles would \emph{slow down} the rise of resistance compared with the model.
\item \textbf{Cost of resistance:} in the model $R$ has no disadvantage without insecticide; in reality \emph{kdr} mosquitoes may be slightly less fit (e.g. slower development) when no insecticide is present, which would also \emph{slow} the spread of $R$ and allow susceptibility to rebound during spray-free periods.
\end{itemize}

\textbf{Task 9 -- Conclusion.} Our model shows that insecticide resistance is not ``created'' by the insecticide: the $R$ allele pre-exists at low frequency, and spraying simply changes the \emph{selective environment}, so that resistant mosquitoes leave more offspring and $\mathrm{freq}(R)$ rises each generation. This is natural selection in action: differential survival and reproduction of heritable variants. Evolution was observed directly, in six ``generations'', in a bag of beads --- and the same process, operating on millions of mosquitoes over many generations per year, explains the pyrethroid resistance now threatening malaria control in Douala and elsewhere in Cameroon. We recommend that the Programme rotate insecticides and combine chemical and environmental control methods, so that no single selection pressure is applied continuously.

\begin{aretenir}[Key points of the investigation]
\begin{itemize}
\item Evolution is a change in \cle{allele frequencies} in a population over generations; natural selection is one of its causes.
\item Selection acts on \emph{phenotypes}, but it is \emph{genotypes} (alleles) that are transmitted.
\item A recessive allele can persist at low frequency by ``hiding'' in heterozygotes.
\item Insecticide resistance is a direct, observable example of natural selection, with major public-health consequences in Cameroon.
\end{itemize}
\end{aretenir}

\begin{attention}[Common mistakes in the report]
\begin{itemize}
\item Forgetting that allele frequencies must always sum to 1 ($p + q = 1$) --- check every row of your table.
\item Dividing by the \emph{original} number of offspring instead of the number of \emph{survivors} when recalculating $p$ after selection.
\item Saying ``the insecticide caused the mutation'': mutations arise at random, before selection; the insecticide only \emph{selects} among existing variants.
\item Confusing the survival of \emph{individuals} with the evolution of the \emph{population}: individuals do not evolve; allele frequencies in populations do.
\end{itemize}
\end{attention}

\newpage

\coursec{Methods and worked examples}

\courssub{Skill 1: Applying the Hardy--Weinberg Principle to a Population}

\begin{methode}[Solving a Hardy--Weinberg Problem]
The Hardy--Weinberg principle states that, in a large, randomly mating population with no mutation, no migration, no selection and no genetic drift, allele and genotype frequencies remain constant from generation to generation. For a gene with two alleles $A$ (frequency $p$) and $a$ (frequency $q$):
\[
p + q = 1 \qquad \text{and} \qquad p^2 + 2pq + q^2 = 1
\]
where $p^2$ = frequency of $AA$, $2pq$ = frequency of $Aa$, $q^2$ = frequency of $aa$.
\begin{enumerate}
\item \textbf{Identify what is given.} Usually the frequency of the \emph{recessive phenotype} (this equals $q^2$) or of the dominant phenotype ($p^2 + 2pq$).
\item \textbf{Extract $q$ (or $p$).} If the recessive phenotype frequency is given, take the square root: $q = \sqrt{q^2}$.
\item \textbf{Find the other allele frequency} using $p = 1 - q$.
\item \textbf{Compute the required genotype frequencies}: $p^2$, $2pq$, $q^2$.
\item \textbf{Convert to numbers of individuals} if asked: multiply the frequency by the population size.
\item \textbf{Check}: the three genotype frequencies must add up to 1 (or the numbers of individuals must add up to the population size).
\end{enumerate}
\textbf{Reflexes:} never take the square root of a \emph{dominant} phenotype frequency; only the recessive phenotype directly gives $q^2$. Carriers (heterozygotes) of a recessive condition are counted by $2pq$, not by $q^2$.
\end{methode}

\begin{exoresolu}[Sickle-cell anaemia in a Cameroonian population]
In a town in the Centre Region of Cameroon, a survey of $10\,000$ newborn babies found that $100$ suffered from sickle-cell anaemia (genotype $Hb^S Hb^S$). Assuming the population is in Hardy--Weinberg equilibrium:
\begin{enumerate}
\item Calculate the frequency of the $Hb^S$ allele and of the normal $Hb^A$ allele.
\item Calculate the number of babies who are carriers (heterozygotes) of the sickle-cell allele.
\item Explain why the $Hb^S$ allele remains common in this region despite being harmful in the homozygous state.
\end{enumerate}
\tcblower
\textbf{Step 1 — From phenotype to allele frequencies.}

The sickle-cell sufferers are homozygous recessive, so their frequency gives $q^2$ directly:
\[
q^2 = \frac{100}{10\,000} = 0.01 \quad\Longrightarrow\quad q = \sqrt{0.01} = 0.1
\]
Then:
\[
p = 1 - q = 1 - 0.1 = 0.9
\]
So the frequency of $Hb^A$ is \textbf{0.9} and that of $Hb^S$ is \textbf{0.1}.

\textbf{Step 2 — Number of carriers.}

The heterozygote frequency is:
\[
2pq = 2 \times 0.9 \times 0.1 = 0.18
\]
Number of carriers:
\[
0.18 \times 10\,000 = \boxed{1\,800 \text{ babies}}
\]
Check: $p^2 = 0.81 \Rightarrow 8\,100$ normal homozygotes; $8\,100 + 1\,800 + 100 = 10\,000$. \checkmark

\textbf{Step 3 — Why the allele persists.}

In regions where malaria is endemic, heterozygotes ($Hb^A Hb^S$) are \emph{more resistant} to severe malaria than normal homozygotes. This is \cle{heterozygote advantage} (balanced polymorphism): natural selection favours the heterozygote, so the $Hb^S$ allele is maintained in the population at a relatively high frequency even though $Hb^S Hb^S$ individuals suffer from a serious disease. This is a classic example of \textbf{balancing selection}.
\end{exoresolu}

\courssub{Skill 2: Interpreting Selection from Population Data}

\begin{methode}[Identifying the Type of Natural Selection]
When given data (table or graph) on a character before and after selection:
\begin{enumerate}
\item \textbf{Draw or visualise the distribution} of the character (e.g. birth mass, beak depth) before selection.
\item \textbf{Mark which phenotypes survive / reproduce} according to the data.
\item \textbf{Compare the mean and the spread} before and after selection.
\item \textbf{Conclude:}
\begin{itemize}
\item Mean \emph{shifts} in one direction, extremes on one side eliminated $\rightarrow$ \cle{directional selection}.
\item Mean \emph{unchanged} but extremes on \emph{both} sides eliminated, spread reduced $\rightarrow$ \cle{stabilising selection}.
\item Mean splits, intermediates eliminated $\rightarrow$ \cle{disruptive selection}.
\end{itemize}
\item \textbf{Justify with the environmental pressure} named in the question (drought, predation, antibiotic, etc.).
\end{enumerate}
\textbf{Reflex:} always quote figures from the data in your justification — examiners award marks for use of data.
\end{methode}

\begin{popfigure}
\begin{center}
\begin{tikzpicture}[scale=0.9]
% Stabilising
\begin{scope}
\draw[->, popline] (0,0) -- (4.2,0) node[below, font=\scriptsize] {phenotype};
\draw[->, popline] (0,0) -- (0,2.2);
\draw[popblue, thick, domain=0.3:3.9, samples=60] plot (\x, {1.8*exp(-((\x-2.1)^2)/0.8)});
\draw[poporange, thick, dashed, domain=0.9:3.3, samples=60] plot (\x, {1.8*exp(-((\x-2.1)^2)/0.25)});
\node[font=\scriptsize\bfseries, popdark] at (2.1,-0.9) {Stabilising};
\end{scope}
% Directional
\begin{scope}[xshift=5.2cm]
\draw[->, popline] (0,0) -- (4.2,0) node[below, font=\scriptsize] {phenotype};
\draw[->, popline] (0,0) -- (0,2.2);
\draw[popblue, thick, domain=0.3:3.9, samples=60] plot (\x, {1.8*exp(-((\x-1.6)^2)/0.6)});
\draw[poporange, thick, dashed, domain=0.3:3.9, samples=60] plot (\x, {1.8*exp(-((\x-2.8)^2)/0.6)});
\draw[->, popdark, thick] (1.6,2.0) -- (2.8,2.0);
\node[font=\scriptsize\bfseries, popdark] at (2.1,-0.9) {Directional};
\end{scope}
% Disruptive
\begin{scope}[xshift=10.4cm]
\draw[->, popline] (0,0) -- (4.2,0) node[below, font=\scriptsize] {phenotype};
\draw[->, popline] (0,0) -- (0,2.2);
\draw[popblue, thick, domain=0.3:3.9, samples=60] plot (\x, {1.8*exp(-((\x-2.1)^2)/0.8)});
\draw[poporange, thick, dashed, domain=0.3:3.9, samples=80] plot (\x, {1.5*exp(-((\x-0.9)^2)/0.2) + 1.5*exp(-((\x-3.3)^2)/0.2)});
\node[font=\scriptsize\bfseries, popdark] at (2.1,-0.9) {Disruptive};
\end{scope}
\end{tikzpicture}
\end{center}
\legende{The three types of natural selection. Solid blue curve: distribution before selection; dashed orange curve: distribution after selection.}
\end{popfigure}

\begin{exoresolu}[Birth mass in human babies]
The table shows the mortality of newborn babies in relation to birth mass in a hospital study.
\begin{center}
\begin{tabularx}{\linewidth}{|l|X|X|X|X|X|}
\hline
\rowcolor{popblueL} Birth mass (kg) & 1.0--2.0 & 2.0--3.0 & 3.0--4.0 & 4.0--5.0 & 5.0--6.0 \\
\hline
Mortality (\%) & 95 & 20 & 4 & 18 & 80 \\
\hline
\end{tabularx}
\end{center}
\begin{enumerate}
\item State the type of natural selection operating and justify your answer using the data.
\item Predict the effect of this selection on the mean birth mass of the population over generations.
\end{enumerate}
\tcblower
\textbf{1. Type of selection.}

Mortality is lowest (4\%) in the intermediate class (3.0--4.0 kg) and rises sharply towards both extremes: 95\% for very light babies (1.0--2.0 kg) and 80\% for very heavy babies (5.0--6.0 kg). Both extreme phenotypes are selected \emph{against}, while the intermediate phenotype is favoured. This is \textbf{stabilising selection}.

\textbf{2. Effect on the population.}

Because babies of intermediate mass survive best and will form the next generation, the \emph{mean} birth mass of the population remains close to the optimum (about 3.5 kg), but the \emph{variance decreases}: the range of birth masses becomes narrower from generation to generation. Stabilising selection therefore \emph{conserves} the existing phenotype rather than changing it — it eliminates deviants and maintains the population well adapted to a stable environment.
\end{exoresolu}

\courssub{Skill 3: Analysing Directional Selection — the Peppered Moth Model}

\begin{methode}[Explaining an Evolutionary Change Step by Step]
For any "how did population X evolve trait Y" question (peppered moth, antibiotic resistance, DDT resistance in mosquitoes), structure the answer in the standard four-point Darwinian logic:
\begin{enumerate}
\item \textbf{Variation:} state that the population shows inherited variation (e.g. some bacteria carry a resistance allele).
\item \textbf{Selection pressure:} name the environmental agent (soot-darkened trees, antibiotic, insecticide) that kills or disadvantages one phenotype.
\item \textbf{Differential survival and reproduction:} the favoured phenotype survives and reproduces more, passing the advantageous allele to offspring.
\item \textbf{Change in allele frequency:} over generations the advantageous allele increases in frequency; the population is said to have \emph{evolved} (evolution = change in allele frequency in a population over time).
\end{enumerate}
\textbf{Reflex:} never write that an individual "developed resistance because it needed it" — that is Lamarckian thinking and loses marks. The variation arises \emph{first} (by mutation), selection acts \emph{afterwards}.
\end{methode}

\begin{exoresolu}[Industrial melanism in \emph{Biston betularia}]
Before 1850 in England, the peppered moth \emph{Biston betularia} existed mainly as a light form, well camouflaged on lichen-covered tree trunks; a rare dark (melanic) form was easily eaten by birds. During industrialisation, soot killed the lichens and darkened the trunks, and the dark form became common, reaching over 90\% in polluted areas. After clean-air laws in the twentieth century, the light form became common again.
\begin{enumerate}
\item Explain, using the principles of natural selection, the rise of the dark form during industrialisation.
\item State the type of selection involved.
\item Explain why the dark form did not disappear completely even before industrialisation.
\end{enumerate}
\tcblower
\textbf{1. Rise of the dark form.}

\begin{itemize}
\item \textbf{Variation:} the moth population contained two inherited colour forms — light and dark — produced originally by mutation.
\item \textbf{Selection pressure:} predation by birds. On soot-blackened trunks the light moths were conspicuous and were eaten preferentially; the dark moths were camouflaged.
\item \textbf{Differential survival:} dark moths survived longer and produced more offspring, transmitting the dark allele.
\item \textbf{Allele frequency change:} generation after generation the dark allele increased in frequency until the dark form exceeded 90\% in industrial areas. The population evolved.
\end{itemize}

\textbf{2. Type of selection.}

This is \textbf{directional selection}: the environmental change shifted the optimum phenotype in one direction (towards dark colour), so the mean of the population moved in that direction.

\textbf{3. Persistence of the dark form before 1850.}

The dark allele was continuously reintroduced by \textbf{mutation}, so even though selection kept it at a very low frequency, it was never completely eliminated. This pre-existing variation is exactly what allowed the population to respond rapidly when the environment changed — natural selection acts on existing variation; it does not create it on demand.
\end{exoresolu}

\begin{savaistu}[Resistance in Cameroonian mosquitoes]
The same Darwinian logic explains the spread of resistance to insecticides such as DDT and pyrethroids in \emph{Anopheles} mosquitoes, the vectors of malaria in Cameroon. A few mosquitoes carry resistance alleles; spraying kills the susceptible ones; the survivors reproduce and the resistance allele rises in frequency. This is why insecticide-treated bed nets must be combined with other control measures and why resistance monitoring is a major public-health activity.
\end{savaistu}

\courssub{Skill 4: Explaining Speciation and Recognising Isolating Mechanisms}

\begin{methode}[Building a Speciation Account]
To explain how one species could split into two (or to classify an isolating mechanism in a given example):
\begin{enumerate}
\item \textbf{Start with one interbreeding population} sharing a common gene pool.
\item \textbf{Interrupt gene flow:} identify the barrier — geographical (mountain, river, sea: \emph{allopatric}) or ecological/behavioural within the same area (\emph{sympatric}).
\item \textbf{Independent evolution:} in isolation, the two sub-populations experience different mutations, different selection pressures and genetic drift, so their allele frequencies diverge.
\item \textbf{Reproductive isolation:} divergence accumulates until, even if the populations meet again, they can no longer interbreed successfully — two \emph{species} now exist.
\item \textbf{Classify the isolating mechanism} if asked:
\begin{itemize}
\item \emph{Pre-zygotic:} geographical, ecological (habitat), seasonal (breeding time), behavioural (courtship), mechanical (incompatible genitalia), gametic (gametes incompatible).
\item \emph{Post-zygotic:} hybrid inviability, hybrid sterility (e.g. the mule), hybrid breakdown.
\end{itemize}
\end{enumerate}
\textbf{Reflex:} the decisive test of speciation is \emph{reproductive isolation}, not visible difference.
\end{methode}

\begin{exoresolu}[Speciation on Mount Cameroon]
A species of sunbird lives in the lowland forests around Mount Cameroon. Following volcanic activity, a lava flow splits the forest, separating the population into an eastern and a western group that cannot cross the barren lava field. After many thousands of years the forest regenerates over the lava and the two groups meet again, but their songs and courtship displays now differ and they no longer interbreed.
\begin{enumerate}
\item Name and describe the type of speciation that occurred.
\item Classify the isolating mechanisms operating (a) during the separation and (b) after the two groups met again.
\item Explain why the two groups diverged genetically during the separation.
\end{enumerate}
\tcblower
\textbf{1. Type of speciation.}

This is \textbf{allopatric speciation}: a \emph{geographical barrier} (the lava field) divided one population into two isolated sub-populations, which then diverged until they became separate species.

\textbf{2. Isolating mechanisms.}
\begin{itemize}
\item (a) During separation: a \textbf{geographical (pre-zygotic) isolating mechanism} — the lava field physically prevented any mating and therefore any gene flow between the groups.
\item (b) After reunion: a \textbf{behavioural (pre-zygotic) isolating mechanism} — the differences in song and courtship displays mean that males and females of the two groups no longer recognise each other as potential mates, so mating does not occur even though they live in the same area.
\end{itemize}

\textbf{3. Why the groups diverged.}

Isolated from each other, the two sub-populations evolved independently:
\begin{itemize}
\item different \textbf{mutations} arose by chance in each group;
\item the two sides of the mountain offered slightly different \textbf{selection pressures} (vegetation, climate, food plants), favouring different alleles;
\item \textbf{genetic drift} caused random changes in allele frequencies, especially in small populations.
\end{itemize}
With no gene flow to mix the two gene pools, these differences accumulated over generations until the populations were genetically so different that they were reproductively isolated — speciation was complete.
\end{exoresolu}

\courssub{Skill 5: Using Evidence to Support Organic Evolution}

\begin{methode}[Answering "How does X provide evidence for evolution?"]
Whether the evidence is fossils, comparative anatomy, embryology, biochemistry or geographical distribution, use the same three-step frame:
\begin{enumerate}
\item \textbf{Describe the observation} precisely (e.g. forelimbs of man, whale, bat and bird share the same bone plan).
\item \textbf{State what it shows} (e.g. these structures are homologous — same basic plan, different functions).
\item \textbf{Link to common descent:} the observation is best explained by inheritance from a common ancestor followed by modification (adaptive radiation / divergent evolution).
\end{enumerate}
\textbf{Key distinctions to memorise:}
\begin{itemize}
\item \cle{Homologous structures}: same origin, different function $\rightarrow$ divergent evolution (pentadactyl limb).
\item \cle{Analogous structures}: different origin, same function $\rightarrow$ convergent evolution (wing of bird vs wing of butterfly).
\item \cle{Vestigial organs}: reduced, functionless remnants (appendix, hind-limb bones of python) $\rightarrow$ inherited from ancestors in which they were functional.
\end{itemize}
\end{methode}

\begin{exoresolu}[Comparative anatomy and biochemistry]
\begin{enumerate}
\item The forelimbs of a man, a bat, a whale and a bird all contain the same set of bones (humerus, radius, ulna, carpals and digits) although they are used for grasping, flying and swimming respectively. Explain how this supports the theory of organic evolution.
\item The wing of a bird and the wing of a butterfly are both used for flight. Why are they \emph{not} evidence of close common ancestry, and what do they illustrate instead?
\item Cytochrome $c$ from humans differs from that of rhesus monkeys by 1 amino acid, from dogs by 11, and from yeast by over 40. What does this pattern show?
\end{enumerate}
\tcblower
\textbf{1. The pentadactyl limb.}

The four forelimbs are built on the same basic plan — one proximal bone (humerus), two distal bones (radius and ulna), then wrist bones and digits — yet perform very different functions. They are \textbf{homologous structures}. The simplest explanation is that man, bat, whale and bird all \emph{inherited this limb plan from a common ancestor}, and that in each lineage the limb was modified by natural selection for a different way of life. This is \textbf{divergent evolution (adaptive radiation)} and is strong evidence of common descent.

\textbf{2. Bird wing versus butterfly wing.}

Although both are flight organs, the bird's wing is a modified bony forelimb covered with feathers, while the butterfly's wing is a chitinous outgrowth of the body wall: their \emph{origin and structure are completely different}. They are \textbf{analogous structures}, produced by \textbf{convergent evolution} — unrelated organisms evolving similar structures because natural selection moulded them for the same function. Analogy therefore does \emph{not} indicate close common ancestry; it shows that similar environmental pressures produce similar adaptations.

\textbf{3. Cytochrome $c$ comparisons.}

The number of amino-acid differences in cytochrome $c$ increases with the evolutionary distance between species: humans and monkeys (very few differences) share a recent common ancestor, while humans and yeast (many differences) diverged in the remote past. This \textbf{biochemical (molecular) evidence} shows that the degree of molecular similarity matches the family tree established from anatomy and fossils — exactly what the theory of common descent predicts.
\end{exoresolu}

\courssub{Skill 6: Comparing Theories of Evolution}

\begin{methode}[Structuring a Comparison of Lamarck, Darwin and Neo-Darwinism]
For an essay asking you to compare or evaluate theories of evolution:
\begin{enumerate}
\item \textbf{State each theory's core claim} in one or two sentences:
\begin{itemize}
\item \emph{Lamarck:} characters acquired during life through use or disuse are inherited ("inheritance of acquired characteristics").
\item \emph{Darwin--Wallace:} inherited variation exists; organisms overproduce; there is a struggle for existence; the best-adapted survive and reproduce (natural selection); favourable traits accumulate over generations.
\item \emph{Neo-Darwinism (modern synthesis):} Darwin's selection combined with Mendelian genetics — mutations and recombination supply variation; evolution is a change in allele frequencies.
\end{itemize}
\item \textbf{Give the decisive objection to Lamarck:} acquired characters (e.g. muscular arms from exercise, a docked tail) are not transmitted to offspring; only changes in the DNA of gametes are heritable.
\item \textbf{Support Darwin with evidence:} industrial melanism, antibiotic resistance, artificial selection in crops and livestock.
\item \textbf{Conclude} that modern evolutionary theory is Darwinian selection operating on genetic variation.
\end{enumerate}
\end{methode}

\begin{exoresolu}[Lamarck versus Darwin — the giraffe's neck]
\begin{enumerate}
\item Contrast the Lamarckian and Darwinian explanations of how the giraffe acquired its long neck.
\item State one observation that disproves the Lamarckian explanation.
\item Explain why Darwin's theory is described as a "theory of natural selection acting on variation".
\end{enumerate}
\tcblower
\textbf{1. The two explanations.}
\begin{itemize}
\item \textbf{Lamarck:} ancestral giraffes had short necks. In reaching for leaves high in trees, each giraffe \emph{stretched} its neck during its lifetime; this acquired lengthening was passed to its offspring, and over many generations of stretching the neck became long. Change is driven by \emph{need and use}, and acquired characters are inherited.
\item \textbf{Darwin:} the ancestral population showed \emph{inherited variation} in neck length — some giraffes were born with slightly longer necks than others. When food was scarce, longer-necked individuals reached more leaves, survived better and produced more offspring, transmitting their long-neck alleles. Generation after generation, natural selection increased the frequency of long-neck alleles, so the average neck length of the population increased.
\end{itemize}

\textbf{2. Disproof of Lamarck.}

Acquired characters are not inherited: for example, the children of a body-builder are not born muscular, and mice whose tails are cut off generation after generation continue to be born with full-length tails. Only changes present in the genes of the gametes can be transmitted.

\textbf{3. Why "selection acting on variation".}

Darwin's mechanism has two inseparable components: (i) \textbf{heritable variation} already present in the population provides the raw material; (ii) \textbf{natural selection} — differential survival and reproduction imposed by the environment — sorts among these variants, increasing the frequency of advantageous alleles. Neither component alone produces evolution: variation without selection changes nothing directionally, and selection without heritable variation has nothing to act upon. Modern genetics (mutation as the source of new alleles, Mendelian inheritance as the transmission mechanism) completed this picture in the neo-Darwinian synthesis.
\end{exoresolu}

\courssub{Skill 7: Testing Evolutionary Data with the Chi-Squared Test}

\begin{methode}[Using $\chi^2$ to Test Whether Observed Ratios Fit a Prediction]
When a genetic cross or population sample gives numbers that "look close" to an expected ratio, test objectively:
\begin{enumerate}
\item \textbf{State the null hypothesis} $H_0$: there is no significant difference between observed and expected numbers (any deviation is due to chance).
\item \textbf{Calculate expected numbers} from the predicted ratio (e.g. $3:1$) using the total sample size.
\item \textbf{Compute} $\displaystyle \chi^2 = \sum \frac{(O - E)^2}{E}$ for each class.
\item \textbf{Find the degrees of freedom}: $df = (\text{number of classes}) - 1$.
\item \textbf{Compare with the critical value} at $p = 0.05$: if $\chi^2 <$ critical value, accept $H_0$ (data fit the prediction); if $\chi^2 >$ critical value, reject $H_0$.
\end{enumerate}
Critical values at $p=0.05$: $df=1 \rightarrow 3.84$; $df=2 \rightarrow 5.99$; $df=3 \rightarrow 7.82$.
\end{methode}

\begin{exoresolu}[Testing a $3:1$ ratio in maize]
A student crosses two heterozygous maize plants ($Aa \times Aa$), where $A$ (purple grain) is dominant to $a$ (yellow grain). Among 400 grains produced, 310 are purple and 90 are yellow. Use the chi-squared test to decide whether these results fit the expected $3:1$ ratio.
\tcblower
\textbf{Step 1 — Null hypothesis.}

$H_0$: the observed numbers do not differ significantly from a $3:1$ ratio; deviations are due to chance.

\textbf{Step 2 — Expected numbers.}

Total $= 400$. Expected purple $= \tfrac{3}{4} \times 400 = 300$; expected yellow $= \tfrac{1}{4} \times 400 = 100$.

\textbf{Step 3 — Compute $\chi^2$.}
\begin{align*}
\chi^2 &= \sum \frac{(O-E)^2}{E} = \frac{(310-300)^2}{300} + \frac{(90-100)^2}{100} \\
&= \frac{100}{300} + \frac{100}{100} = 0.33 + 1.00 = 1.33
\end{align*}

\textbf{Step 4 — Degrees of freedom.}

There are 2 classes (purple, yellow), so $df = 2 - 1 = 1$.

\textbf{Step 5 — Compare and conclude.}

Critical value at $p = 0.05$ for $df = 1$ is $3.84$. Since $\chi^2 = 1.33 < 3.84$, we \textbf{accept the null hypothesis}: the difference between observed (310:90) and expected (300:100) numbers is not statistically significant and can be attributed to chance. The data are therefore \textbf{consistent with monohybrid inheritance} giving a $3:1$ phenotypic ratio — the cross behaves as Mendelian genetics predicts.

\textbf{Examiner's note:} a very common error is to compare $\chi^2$ with the wrong critical value (always check the degrees of freedom) or to conclude that $H_0$ is "proved" — we only say the data are \emph{consistent with} it.
\end{exoresolu}

\begin{attention}[Frequent Errors in GCE Answers on Evolution]
\begin{itemize}
\item Writing that organisms "evolve because they need to" or that an individual "adapts itself" — evolution acts on \emph{populations}, through selection of \emph{pre-existing} heritable variation.
\item Confusing \emph{homologous} with \emph{analogous} structures — always mention both origin and function.
\item Taking the square root of the dominant phenotype frequency in Hardy--Weinberg problems — only $q^2$ (recessive phenotype) can be square-rooted directly.
\item Saying the mule proves horses and donkeys are one species — hybrid \emph{sterility} is a post-zygotic isolating mechanism proving they are \emph{different} species.
\item Forgetting that natural selection \emph{preserves} favourable traits but does not \emph{create} them; mutations create new alleles.
\end{itemize}
\end{attention}

\newpage

\coursec{Exercises}

\courssub{I check my knowledge}

\begin{exercice}[MCQ: Paper 1 style MCQ]
\begin{enumerate}
\item Which of the following is the unit of natural selection?
\begin{enumerate}[label=\alph*)]
\item the individual
\item the population
\item the gene
\item the ecosystem
\end{enumerate}
\item In a randomly mating population of $N = 100$ individuals, $4\,\%$ of the population is affected by sickle cell anaemia (genotype $\mathrm{HbS}/\mathrm{HbS}$). Assuming Hardy--Weinberg equilibrium, the frequency of the $\mathrm{HbS}$ allele is:
\begin{enumerate}[label=\alph*)]
\item $0.04$
\item $0.2$
\item $0.32$
\item $0.8$
\end{enumerate}
\end{enumerate}
\end{exercice}

\begin{attention}[Common mistake]
Do not confuse the \emph{genotypic} frequency $q^{2}$ (proportion of affected $\mathrm{HbS}/\mathrm{HbS}$ individuals) with the \emph{allelic} frequency $q$. You must take the square root: $q = \sqrt{q^{2}}$. Choosing answer (a) $0.04$ is the classic trap in Paper 1.
\end{attention}

\newpage

\coursec{Solutions to the exercises}

\begin{corrige}[Exercise 1 — MCQ: Paper 1 style MCQ]
\begin{enumerate}
\item \textbf{Answer: (b) the population.}

Natural selection \emph{acts} on the phenotype of individual organisms: within a population, some individuals survive and reproduce more successfully than others because of their heritable characteristics. However, the result of selection — a change in allele frequencies from one generation to the next, i.e. \cle{evolution} — can only be measured at the level of the \cle{population}. An individual does not evolve during its lifetime; its genotype is fixed at fertilisation. Hence the classic formulation: \emph{the individual is selected, but the population evolves}. Options (a), (c) and (d) are therefore incorrect: the individual is the target of selection but not its evolutionary unit, the gene is the unit of inheritance, and the ecosystem is far too broad a level.

\item \textbf{Answer: (b) $0.2$.}

Sickle cell anaemia is expressed only in homozygous recessive individuals ($\mathrm{Hb^S}/\mathrm{Hb^S}$). Under Hardy--Weinberg equilibrium, the genotypic frequencies are
\[
p^{2} \;(\mathrm{Hb^A}/\mathrm{Hb^A}) + 2pq \;(\mathrm{Hb^A}/\mathrm{Hb^S}) + q^{2} \;(\mathrm{Hb^S}/\mathrm{Hb^S}) = 1,
\]
where $p$ is the frequency of the normal allele $\mathrm{Hb^A}$ and $q$ that of the sickle allele $\mathrm{Hb^S}$.

The affected individuals represent $4\,\%$ of the population, so
\[
q^{2} = 0.04 \quad\Longrightarrow\quad q = \sqrt{0.04} = 0.2.
\]
The frequency of the $\mathrm{Hb^S}$ allele is therefore $q = 0.2$, and that of the normal allele is $p = 1 - q = 0.8$.

\textbf{Check (good examination habit):} carrier frequency $= 2pq = 2 \times 0.8 \times 0.2 = 0.32$, i.e. $32\,\%$ of the population are healthy carriers ($\mathrm{Hb^A}/\mathrm{Hb^S}$) — a plausible value for a malaria-exposed population such as that of Cameroon, where the heterozygote advantage maintains the $\mathrm{Hb^S}$ allele. Note that option (c) $0.32$ is precisely this carrier frequency, included as a distractor.
\end{enumerate}

\textbf{Examiner's expectations (Paper 1):} for this type of item, the candidate must (i) identify the affected class as the homozygous recessive class $q^{2}$, (ii) take the square root to obtain the allele frequency, and (iii) use $p + q = 1$ if the other allele is required. No working is shown on the answer sheet, so accuracy of the mental steps is essential.
\end{corrige}

\begin{attention}[Common mistake]
Do not confuse the \emph{genotypic} frequency $q^{2}$ (proportion of affected $\mathrm{Hb^S}/\mathrm{Hb^S}$ individuals) with the \emph{allelic} frequency $q$. You must take the square root: $q = \sqrt{q^{2}}$. Choosing answer (a) $0.04$ is the classic trap in Paper 1. Equally, $0.32$ is the frequency of \emph{carriers} ($2pq$), not of the allele — read the question twice before shading your answer.
\end{attention}

\newpage

\coursec{Summary sheet}

\begin{popfigure}
\begin{tikzpicture}[
  grow cyclic,
  mindmap,
  concept color=popblue,
  every node/.style={concept, text=white, font=\small},
  level 1 concept/.append style={sibling angle=72, level distance=3.4cm, font=\small},
  level 2 concept/.append style={sibling angle=45, level distance=2.1cm, font=\scriptsize}
]
\node[concept, concept color=popdark, font=\bfseries\small] {Natural Selection \& Evolution}
  child[concept color=popblue] { node {Variation \& Populations}
    child { node {Gene pool, allele frequency} }
    child { node {Hardy--Weinberg $p^2+2pq+q^2=1$} }
    child { node {Mutation, drift, gene flow} }
  }
  child[concept color=popgreen] { node {Natural Selection}
    child { node {Stabilising} }
    child { node {Directional} }
    child { node {Disruptive} }
  }
  child[concept color=poporange] { node {Speciation}
    child { node {Allopatric / sympatric} }
    child { node {Isolating mechanisms} }
  }
  child[concept color=poppurple] { node {Evidence of Evolution}
    child { node {Fossils} }
    child { node {Comparative anatomy, embryology} }
    child { node {Molecular evidence} }
  }
  child[concept color=poppink] { node {Theories}
    child { node {Lamarck} }
    child { node {Darwin--Wallace} }
    child { node {Neo-Darwinism} }
  };
\end{tikzpicture}
\legende{Mind map of Chapter BIO09: Natural selection and evolution.}
\end{popfigure}

\begin{bilan}
\textbf{Key definitions}\par
\begin{itemize}
\item \cle{Evolution}: the gradual, cumulative change in the heritable characteristics of populations over successive generations.
\item \cle{Population}: a group of organisms of the same species living in the same area at the same time, capable of interbreeding.
\item \cle{Gene pool}: the total of all alleles of all genes in a population.
\item \cle{Allele frequency}: the proportion of a given allele among all alleles of that gene in the population.
\item \cle{Natural selection}: the differential survival and reproduction of individuals owing to differences in phenotype; the best-adapted leave more offspring, so favourable alleles increase in frequency.
\item \cle{Fitness}: the relative contribution an individual makes to the gene pool of the next generation.
\item \cle{Species}: a group of organisms that can interbreed in nature to produce fertile offspring.
\item \cle{Speciation}: the formation of new species from an ancestral population.
\item \cle{Adaptive radiation}: the rapid diversification of one ancestral stock into many species adapted to different niches (e.g.\ Darwin's finches).
\end{itemize}

\textbf{Laws and formulas to know}\par
\begin{itemize}
\item \cle{Hardy--Weinberg principle}: in a large, randomly mating population with no mutation, migration, selection or drift, allele and genotype frequencies remain constant from generation to generation.
\item With $p$ = frequency of dominant allele, $q$ = frequency of recessive allele:
\[ p + q = 1 \qquad\text{and}\qquad p^2 + 2pq + q^2 = 1 \]
where $p^2$ = frequency of homozygous dominant, $2pq$ = heterozygotes, $q^2$ = homozygous recessive.
\item Conditions of validity: large population, random mating, no mutation, no migration, no selection.
\item \cle{Genetic drift}: random change of allele frequencies, strongest in small populations (founder effect, bottleneck effect).
\end{itemize}

\textbf{Types of natural selection}\par
\begin{itemize}
\item \cle{Stabilising selection}: favours the average phenotype, eliminates extremes (e.g.\ human birth mass).
\item \cle{Directional selection}: favours one extreme; the mean shifts (e.g.\ industrial melanism in \emph{Biston betularia}, antibiotic resistance in bacteria).
\item \cle{Disruptive selection}: favours both extremes against the intermediate; may lead to polymorphism and speciation.
\end{itemize}

\textbf{Speciation and isolating mechanisms}\par
\begin{itemize}
\item \cle{Allopatric speciation}: populations separated by a geographical barrier diverge until they can no longer interbreed.
\item \cle{Sympatric speciation}: speciation without geographical separation (e.g.\ polyploidy in plants).
\item \cle{Pre-zygotic barriers}: habitat, seasonal (temporal), behavioural, mechanical isolation; gamete incompatibility.
\item \cle{Post-zygotic barriers}: hybrid inviability, hybrid sterility (e.g.\ the mule), hybrid breakdown.
\end{itemize}

\textbf{Evidence of organic evolution}\par
\begin{itemize}
\item \cle{Palaeontology}: fossils show a sequence of forms through geological time; transitional forms (e.g.\ \emph{Archaeopteryx}).
\item \cle{Comparative anatomy}: homologous structures (pentadactyl limb) indicate common ancestry; analogous structures show convergent evolution; vestigial organs.
\item \cle{Comparative embryology}: similar early embryos of vertebrates.
\item \cle{Molecular evidence}: DNA and protein sequence comparisons; the more similar the sequences, the closer the relationship.
\item \cle{Geographical distribution} (biogeography) and observed selection (resistance in mosquitoes and bacteria in Cameroon).
\end{itemize}

\textbf{Theories of evolution}\par
\begin{itemize}
\item \cle{Lamarckism}: inheritance of acquired characteristics through use and disuse — now rejected.
\item \cle{Darwin--Wallace theory}: overproduction, variation, struggle for existence, survival of the fittest, descent with modification.
\item \cle{Neo-Darwinism (modern synthesis)}: Darwinian selection combined with Mendelian genetics; mutations supply variation, selection acts on phenotypes, evolution is change in allele frequencies.
\end{itemize}

\textbf{Methods}\par
\begin{itemize}
\item To use Hardy--Weinberg: (1) compute $q^2$ from the proportion of recessive phenotype; (2) take $q=\sqrt{q^2}$; (3) $p=1-q$; (4) deduce $p^2$ and $2pq$.
\item To explain selection: variation $\rightarrow$ selection pressure $\rightarrow$ differential survival/reproduction $\rightarrow$ allele frequency change over generations.
\end{itemize}

\textbf{Common mistakes}\par
\begin{itemize}
\item Saying individuals evolve: \emph{populations} evolve, individuals are selected.
\item Confusing homologous (common origin) with analogous (same function) structures.
\item Believing mutations appear ``because they are needed'': mutations are random; selection is not.
\item Forgetting to take the square root of $q^2$ before computing $p$.
\item Attributing Lamarck's ideas to Darwin.
\end{itemize}

\textbf{GCE tips}\par
\begin{itemize}
\item Always define key terms (species, natural selection, gene pool) before developing an essay.
\item Use named examples: peppered moth, sickle cell and malaria (heterozygote advantage in Cameroon), Darwin's finches.
\item In data questions, show every step of Hardy--Weinberg calculations with correct notation.
\item For evidence questions, give one example per line of evidence and state what it shows.
\end{itemize}
\end{bilan}

\findecours

\end{document}
